% !TEX root = ../main.tex
\chapter{Portfolio Theory and Portfolio Management}\label{ch:portfolio}
\begin{paracol}{2}

\begin{sources}
\begin{tabularx}{\linewidth}{@{}l l L@{}}
\toprule
\textbf{Lecture} & \textbf{Speaker} & \textbf{Content}\\
\midrule
2013 L14 & P.~Kempthorne & Markowitz mean--variance analysis (two assets, then $m$ assets with Lagrange multipliers); the risk-free asset, the market portfolio, Tobin's separation theorem and the capital market line; von Neumann--Morgenstern utility; portfolio constraints (long-only, holding, turnover, benchmark, tracking error, factor, integer); estimation of means and covariances; alternative risk measures (MAD, semi-variance, VaR, CVaR, coherent measures, skewness and kurtosis). Live demonstration of Case Study~6 (simulation and sector ETFs).\\
2013 L16 & J.~Xia & Blackboard lecture: goals of portfolio management, the return--volatility map of asset classes, portfolio theory through special cases of two assets, Sharpe ratio, beta, the Kelly formula, the 60/40 portfolio, risk parity and leverage, the ``free lunch'' of rebalancing, crowd synchronisation (Millennium Bridge, metronomes) and the limits of the theory. Long question session.\\
2024 L13 & J.~Xia & The same programme in a slide version with new material from the speaker's own research: the endowment model, the loss of meaning of volatility as risk, the gain--loss ratio $(G-L)/(G+L)$ for sizing, a feedback model of crowd behaviour and the origin of power laws. Questions on taxes, expected versus worst loss, stop losses.\\
\bottomrule
\end{tabularx}

\smallskip
\textbf{Files.} 2024 slides \file{mit18_642_f24_lec13.pdf} (38 pages). 2013 slides \file{MIT18_S096F13_lecnote14.pdf}
(34 pages) and \file{lecnote16.pdf} (2 pages: only the outline, the lecture was given on the blackboard);
2013 Case Study~6 \file{CaseStudy6.pdf} (portfolio theory with US sector ETFs, 14 pages); the R output
\file{Smltn_TwoAst.pdf} (simulation of two-asset portfolios, 6 pages) and the four sector-ETF output files
\file{ETF_pridA_15.pdf}, \file{ETF_pridA_30.pdf} (2009--2013, maximum allocation $15\%$ and $30\%$),
\file{ETF_pridB_15.pdf}, \file{ETF_pridB_30.pdf} (2003--2006); the names ``periodA'' and ``periodB'' are those of
the case study text.
\textbf{Problem sets.} 2013 PS6 Problems~3--5 (two-asset and $m$-asset minimum-variance portfolios,
equicorrelation); Problems~1--2 of the same set are time series (\cref{ch:timeseries}). The exercises are at the
end of the chapter.
\textbf{Not covered or not available.} The R code that produced the case-study output is not in the download (only
its printed output), so the numbers of \cref{ch10:sec-case} are quoted, and checked by an independent computation
from the printed input tables. The video shown by the speaker in both years and the 2013 blackboard are not in the
notes; \file{lecnote16} contains only the outline of the lecture. The second half of 2013 L14 mentions ``sector
market-neutral models'' with no data (\cref{ch10:sec-case}).
\end{sources}

\section*{Motivation}
Every investor faces the same question: given a list of things one could hold, how much of each? The 2013 and 2024
lectures answer it twice. Kempthorne (2013 L14) gives the mathematical theory that earned Harry Markowitz a Nobel
prize: describe a portfolio by its expected return and its variance, and find the portfolios that offer the
largest return for a given variance. The theory is a piece of applied linear algebra (\cref{ch:linalg}): a
quadratic form to be minimised under linear constraints, solved with Lagrange multipliers, whose solution is
governed by the inverse covariance matrix $\Sigma^{-1}$. Xia (2013 L16, 2024 L13), who has managed money for
a bank and for the Harvard endowment, then asks what survives in practice: the inputs are estimated with large
errors, volatility is a poor measure of risk, and markets are populated by agents who react to each other.

The chapter follows this arc.
\begin{itemize}
  \item \cref{ch10:sec-problem}: the objectives, constraints and horizon that define a portfolio problem.
  \item \cref{ch10:sec-two,ch10:sec-mvo,ch10:sec-rf}: Markowitz theory for two assets, for $m$ assets, and with a
        riskless asset (efficient frontier, tangency portfolio, capital market line), with complete derivations;
        \cref{ch10:sec-capm} links it to beta and the CAPM (\cref{ch03:thm-capm}).
  \item \cref{ch10:sec-rebal,ch10:sec-rp}: diversification, rebalancing, leverage and risk parity.
  \item \cref{ch10:sec-utility,ch10:sec-kelly}: utility theory as the foundation of mean--variance analysis, and
        sizing by the Kelly criterion and by Xia's gain--loss ratio.
  \item \cref{ch10:sec-constr,ch10:sec-est,ch10:sec-risk}: constraints, estimation error, and other risk measures.
  \item \cref{ch10:sec-crowd}: the ``beyond Markowitz'' ideas of the 2024 lecture (crowds, power laws).
  \item \cref{ch10:sec-case}: the 2013 case study and the notes on sources and errata.
\end{itemize}
\emph{Notation.} The slides call the vector of expected returns $\bm\alpha$; here it is $\bm\mu$, to keep
$\alpha$ for the regression intercept of \cref{def:alphabeta}. Vectors are columns, $\bm 1=(1,\dots,1)\T$,
$\Sigma$ is the (positive definite) covariance matrix of the returns $\bm R$ of $m$ risky assets and
$\sigma_i^2=\Sigma_{ii}$.

% ==================================================================
\section{The portfolio construction problem}\label{ch10:sec-problem}

Both Xia lectures start with the same exercise: each student writes on a sheet the composition of a
hypothetical portfolio (\$10{,}000 in 2024; in 2013 anything from a \$1{,}000 saving to a fund's start-up capital), in percent, in five minutes and
without being told the goal \ts{24}{13}{3:14}\ts{13}{16}{3:10}. The answers become data for the rest of the hour.
In 2024 they were mostly index funds (S\&P~500, QQQ), Treasuries and cash, one option strategy, no cryptocurrency
(in 2021 the class had picked Dogecoin) \ts{24}{13}{9:26}\ts{24}{13}{11:35}; in 2013 they ranged from bonds and
S\&P index funds to timberland, rare coins and lottery tickets \ts{13}{16}{10:16}\ts{13}{16}{11:16}.

\subsection{Objective, loss tolerance, horizon}
Before choosing what to hold one must say what the portfolio is \emph{for}. The considerations listed by Xia are:
\begin{itemize}
  \item a \emph{return objective} (10\%, 50\%, 100\%?) and a \emph{time horizon}, which he deliberately left out
        of the exercise \ts{24}{13}{4:15}\ts{24}{13}{18:01};
  \item a \emph{loss tolerance}: how much of \$10{,}000 can be lost before one stops, and drawdown limits
        \ts{24}{13}{15:54};
  \item an \emph{edge}: why should this investment beat the alternatives \ts{24}{13}{5:17};
  \item \emph{constraints}: liquidity and spending needs, exclusions (fossil fuels, tobacco), inflation, absolute
        versus relative returns against a benchmark (the excess return being the \emph{alpha} of
        \cref{def:alphabeta}) \ts{24}{13}{14:48}\ts{24}{13}{16:56};
  \item \emph{asset--liability matching}: after 2008 many asset owners learned that when the portfolio draws down
        the liabilities dominate the balance sheet and spending must be cut \ts{24}{13}{18:01};
  \item \emph{peer pressure and career risk}: managers are compared quarter to quarter although the owner's
        horizon is long \ts{24}{13}{19:01}.
\end{itemize}
The 2013 version adds the individual's life cycle: income peaks around age 50, spending follows a different
curve, and the portfolio must bridge the gap; the same question is asked of a pension fund, an endowment and a
trading desk (how much capital to allocate to each trader, or to each of 30 strategies in a new hedge fund)
\ts{13}{16}{14:32}\ts{13}{16}{17:43}.
In 2024 the objective is written (slide~7) as $\max R/\sigma$ subject to constraints \ts{24}{13}{14:48}.

\subsection{Layers of decisions and the risk--return map}
The slides separate three layers \ts{24}{13}{7:22}: choice of \emph{markets and instruments};
\emph{data, signals, models, factors, predictions and strategies}; and finally \emph{sizing, capital allocation,
portfolio construction, optimisation and risk management}. Only the last is the topic of this chapter. It comes
down to comparing investments through their return and their risk, first taken to be the volatility
\ts{24}{13}{8:25}. On the map with volatility on the horizontal axis (2024 slide~6, 2013 blackboard),
cash sits on the vertical axis with zero volatility and a positive return (over 4\% in October 2024,
\ts{24}{13}{11:35}); bonds have a little more of both, stocks more, an index fund less volatility than a single stock,
commodities more, private equity and venture capital (illiquid, with a wide range of outcomes) are far up to the
right \ts{24}{13}{12:39}\ts{24}{13}{13:44}\ts{13}{16}{24:06}. Two extremes anchor the picture: a lottery ticket has a
return near $-100\%$ and a small volatility, and the fair coin flip has zero return and $100\%$ volatility
\ts{13}{16}{23:02}. A rational investor discards every point that has another point above and to the left of
it, which is the idea of the efficient frontier \ts{13}{16}{26:13}\ts{13}{16}{27:18}.

\begin{finance}[The endowment model]
A university endowment is a perpetual portfolio. It spends about $5\%$ of its value each year, and inflation is
about $3\%$, so the nominal return target is about $8\%$ (more precisely $1.05\times1.03-1\approx8.2\%$),
over horizons ``longer than ten years'' \ts{24}{13}{21:07}\ts{24}{13}{22:12}. The slide adds that liquidity
must cover $30$--$40\%$ of the annual operating budget, and Xia says that roughly $40\%$ of a university's budget is now
paid from endowment spending, as research grants shrink \ts{24}{13}{22:12}. An $8\%$ target is not easy to obtain
consistently when bonds yield little, so endowments hold a long list of strategies \ts{24}{13}{23:15}: cash,
government and corporate bonds and credit, macro and relative-value hedge funds, quantitative funds (trend
following/CTA, statistical arbitrage), long/short equity, multi-manager platforms, private equity, real estate and
natural resources, and newer assets (crypto, intellectual property, legal claims). The ``endowment model'' is
characterised by external rather than internal managers, public and private markets, generalists and specialists,
absolute and relative returns, mostly active management with some passive beta, and peer comparison
\ts{24}{13}{25:20}\ts{24}{13}{26:27}. Its central skill is manager selection (``can you predict fund performance,
luck or skill?''). Endowments are tax-exempt, so taxes are ignored; a family office would tilt to
private equity to defer realised gains \ts{24}{13}{1:13:53}\ts{24}{13}{1:14:54}.
\end{finance}

\subsection{The classic problem}
Assume the objectives and the loss tolerance are given and that one can predict the return, volatility and
correlation matrix of every investment. What fraction should go into each? The classic problem is to maximise the
portfolio return and minimise its risk \ts{24}{13}{27:30}. Three observations of the lecture frame the rest of
the chapter \ts{24}{13}{28:30}\ts{24}{13}{29:35}. First, the problem is one of \emph{sizing} at every level: for asset
classes it is asset allocation, and because ten to fifteen classes make optimisation hard, practitioners often
reduce them to three to five \emph{risk factors} (equity, bonds, inflation, currency, credit; see
\cref{ch:pca}). Second, risk management is the same problem seen from the other side, with the added concerns of
concentration, illiquidity and unwanted risks. Third, everything depends on predictions (``capital market
assumptions'') which are unreliable (\cref{ch10:sec-est}).

% ==================================================================
\section{Two assets: diversification and the feasible set}\label{ch10:sec-two}

Consider two risky assets with returns $R_1,R_2$, expected returns $\mu_1,\mu_2$, volatilities
$\sigma_1,\sigma_2$ and correlation $\rho$, and the fully invested portfolios with weights $(1-w,w)$, $0\le w\le1$
(no short sales) \ts{13}{14}{5:21}\ts{24}{13}{30:39}:
\begin{gather}\label{ch10:eq-two}
  R_w=(1-w)R_1+wR_2,\qquad \mu_w=(1-w)\mu_1+w\mu_2,\notag\\
  \sigma_w^2=(1-w)^2\sigma_1^2+w^2\sigma_2^2+2w(1-w)\rho\sigma_1\sigma_2 .
\end{gather}
The \emph{feasible set} is the curve $\{(\sigma_w,\mu_w):0\le w\le1\}$ in the volatility--return plane
\ts{13}{14}{7:28}. Since $\mu_w$ is affine in $w$ and $\sigma_w^2$ is quadratic, $\sigma_w^2$ is a quadratic
function of $\mu_w$: the curve is one branch of a hyperbola (a parabola in the $(\mu,\sigma^2)$ plane)
\ts{13}{14}{27:32}. Its shape is governed by $\rho$, illustrated by the special cases of both Xia lectures.

\begin{itemize}
  \item $\rho=1$: $\sigma_w=(1-w)\sigma_1+w\sigma_2$ is linear in $w$, hence so is $\mu$ as a function of $\sigma$:
        a straight segment between the two assets and no diversification gain
        \ts{24}{13}{31:39}\ts{13}{16}{34:34}.
  \item $\rho=-1$: $\sigma_w=|(1-w)\sigma_1-w\sigma_2|$, two segments meeting on the return axis at
        $w=\sigma_1/(\sigma_1+\sigma_2)$, where the portfolio has \emph{zero} variance
        \ts{24}{13}{32:39}\ts{13}{14}{17:06}. A portfolio without risk must earn the riskless rate, otherwise
        there is arbitrage; hedges of a security by a future exploit exactly this \ts{13}{14}{18:06}\ts{13}{14}{19:06}.
  \item $\rho=0$: the curve bends to the left of the straight line, and the variance cannot be reduced to zero.
  \item $\sigma_1=0$ (asset~1 is cash with return $r_0$): $\sigma_w=w\sigma_2$ and
        $\mu_w=r_0+(\mu_2-r_0)\sigma_w/\sigma_2$, a straight line with slope $(\mu_2-r_0)/\sigma_2$, the
        \emph{capital allocation line} of \cref{ch10:sec-rf} \ts{24}{13}{33:41}\ts{24}{13}{34:45}\ts{13}{16}{35:37}.
\end{itemize}

\begin{proposition}[Minimum-variance portfolio of two assets]\label{ch10:prop-two}
Let $\sigma_1^2+\sigma_2^2-2\rho\sigma_1\sigma_2>0$ (that is, not $\rho=1$ with $\sigma_1=\sigma_2$). The variance
\eqref{ch10:eq-two} is minimised over $w\in\R$ at
\[
  w^*=\frac{\sigma_1^2-\rho\sigma_1\sigma_2}{\sigma_1^2+\sigma_2^2-2\rho\sigma_1\sigma_2},\qquad
  \sigma^{*2}=\frac{\sigma_1^2\sigma_2^2(1-\rho^2)}{\sigma_1^2+\sigma_2^2-2\rho\sigma_1\sigma_2}.
\]
The weight $w^*$ lies in $[0,1]$ (the minimum is attained by a long-only portfolio) iff
$\rho\le\min(\sigma_1/\sigma_2,\sigma_2/\sigma_1)$; otherwise the minimum over $[0,1]$ is at the endpoint with smaller
volatility. If $\sigma_1=\sigma_2=\sigma$ then $w^*=\tfrac12$ and $\sigma^{*2}=\sigma^2(1+\rho)/2$.
\end{proposition}
\begin{proof}
Differentiating \eqref{ch10:eq-two}: $\tfrac12\,\dd\sigma_w^2/\dd w=-(1-w)\sigma_1^2+w\sigma_2^2+(1-2w)\rho\sigma_1\sigma_2$,
which vanishes at $w^*$; the second derivative $\sigma_1^2+\sigma_2^2-2\rho\sigma_1\sigma_2$ is positive, so this is the
global minimum. Substituting gives $\sigma^{*2}$ (use $1-w^*=(\sigma_2^2-\rho\sigma_1\sigma_2)/(\dots)$).
The numerator of $w^*$ is $\ge0$ iff $\rho\le\sigma_1/\sigma_2$, and that of $1-w^*$ is $\ge0$ iff
$\rho\le\sigma_2/\sigma_1$.
\end{proof}
For $\rho=0$ the formula gives $w^*=\dfrac{1/\sigma_2^2}{1/\sigma_1^2+1/\sigma_2^2}$: the assets are weighted
\emph{inversely to their variances}, as computed by Kempthorne on the board \ts{13}{14}{12:52}\ts{13}{14}{13:54}
(Exercise~\ref{ch10:ex-ps6-3} extends it).
Two consequences of \cref{ch10:prop-two} explain the pictures of the lectures.
\begin{itemize}
  \item Adding a little of the riskier
      asset~2 to asset~1 \emph{reduces} the variance as long as $\rho<\sigma_1/\sigma_2$
      (the derivative of $\sigma_w^2$ at $w=0$ is $2\sigma_1(\rho\sigma_2-\sigma_1)$, negative iff $\rho<\sigma_1/\sigma_2$). Even for $\rho=0$ diversification lowers the risk and, when asset~2 has the higher return,
      also \emph{raises} the return: every portfolio to the left of the minimum-variance point is dominated
      \ts{13}{14}{14:57}\ts{13}{14}{15:59}. For $\rho$ larger than the ratio of volatilities the minimum is the safer asset alone.
  \item The equal-volatility, zero-correlation case gives $\sigma^*=\sigma/\sqrt2$: ``a significant reduction of
      the risk by half-and-half'' \ts{13}{16}{31:31}\ts{13}{16}{33:33}; for three such assets it is $\sigma/\sqrt3$
      \ts{13}{16}{41:28}, and for $m$ assets $\sigma/\sqrt m$ (Exercise~\ref{ch10:ex-ps6-4}).
\end{itemize}

\begin{example}[The two assets of the 2013 simulation]\label{ch10:ex-two}
Take $\mu=(0.15,0.25)$, $\sigma=(0.20,0.30)$, the values printed in the plots of Case Study~6
(\cref{ch10:sec-case} explains why these are not the values read out in class). Then $\sigma_1/\sigma_2=2/3$ and
\begin{center}
\begin{tabular}{@{}lrrrrrr@{}}
\toprule
$\rho$ & $-1$ & $-0.4$ & $0$ & $0.4$ & $0.8$ & $1$\\
\midrule
$w^*$ (asset~2) & 0.400 & 0.360 & 0.308 & 0.195 & $-0.235$ & $-2$\\
$\sigma^*$ (long-only) & 0 & 0.130 & 0.166 & 0.192 & 0.200 & 0.200\\
$\mu^*$ (long-only) & 0.190 & 0.186 & 0.181 & 0.170 & 0.150 & 0.150\\
\bottomrule
\end{tabular}
\end{center}
For $\rho=0.8>2/3$ and $\rho=1$ the unconstrained $w^*$ is negative (a short position in asset~2), and the
long-only minimum is asset~1 alone. This is why in the last panel of \file{Smltn_TwoAst} (correlation $0.8$) the
minimum-variance portfolio (cyan curve) coincides with asset~1. All values were computed from
\cref{ch10:prop-two}.
\end{example}

\begin{figure}[ht]
\centering
\begin{tikzpicture}
\begin{axis}[width=0.86\linewidth,height=6.6cm,xlabel={volatility $\sigma_w$},ylabel={expected return $\mu_w$},
  xmin=-0.005,xmax=0.36,ymin=0.135,ymax=0.265,xtick={0,0.1,0.2,0.3},scaled ticks=false,/pgf/number format/fixed,grid=major,grid style={black!10},
  legend style={font=\footnotesize,at={(0.5,1.03)},anchor=south,legend columns=5},
  tick label style={font=\small},label style={font=\small}]
  \addplot[thick,thmcol] coordinates {(0.2,0.15) (0,0.19) (0.3,0.25)};
  \addplot[thick,cscol] coordinates {(0.2000,0.1500) (0.1845,0.1550) (0.1702,0.1600) (0.1575,0.1650) (0.1467,0.1700) (0.1383,0.1750) (0.1327,0.1800) (0.1304,0.1850) (0.1315,0.1900) (0.1358,0.1950) (0.1432,0.2000) (0.1531,0.2050) (0.1652,0.2100) (0.1789,0.2150) (0.1940,0.2200) (0.2101,0.2250) (0.2270,0.2300) (0.2446,0.2350) (0.2626,0.2400) (0.2811,0.2450) (0.3000,0.2500)};
  \addplot[thick,intcol] coordinates {(0.2000,0.1500) (0.1906,0.1550) (0.1825,0.1600) (0.1759,0.1650) (0.1709,0.1700) (0.1677,0.1750) (0.1664,0.1800) (0.1671,0.1850) (0.1697,0.1900) (0.1741,0.1950) (0.1803,0.2000) (0.1879,0.2050) (0.1970,0.2100) (0.2072,0.2150) (0.2184,0.2200) (0.2305,0.2250) (0.2433,0.2300) (0.2568,0.2350) (0.2707,0.2400) (0.2852,0.2450) (0.3000,0.2500)};
  \addplot[thick,fincol] coordinates {(0.2000,0.1500) (0.1965,0.1550) (0.1940,0.1600) (0.1925,0.1650) (0.1920,0.1700) (0.1927,0.1750) (0.1944,0.1800) (0.1971,0.1850) (0.2008,0.1900) (0.2054,0.1950) (0.2110,0.2000) (0.2173,0.2050) (0.2243,0.2100) (0.2320,0.2150) (0.2404,0.2200) (0.2492,0.2250) (0.2586,0.2300) (0.2684,0.2350) (0.2786,0.2400) (0.2891,0.2450) (0.3000,0.2500)};
  \addplot[thick,defcol] coordinates {(0.2,0.15) (0.3,0.25)};
  \legend{$\rho=-1$,$\rho=-0.4$,$\rho=0$,$\rho=0.4$,$\rho=1$}
  \addplot[only marks,mark=*,mark size=2pt,black,forget plot] coordinates {(0.2,0.15) (0.3,0.25)};
  \node[font=\footnotesize,anchor=east] at (axis cs:0.19,0.150) {asset 1};
  \node[font=\footnotesize,anchor=west] at (axis cs:0.305,0.25) {asset 2};
  \addplot[only marks,mark=square*,mark size=2pt,intcol,forget plot] coordinates {(0.1664,0.1808)};
  \draw[->,intcol] (axis cs:0.235,0.163) -- (axis cs:0.170,0.1790);
  \node[font=\footnotesize,anchor=west,intcol] at (axis cs:0.235,0.160) {min.\ variance, $\rho=0$};
\end{axis}
\end{tikzpicture}
\caption{Feasible set of the portfolios $(1-w)R_1+wR_2$, $0\le w\le1$, for $\mu=(0.15,0.25)$, $\sigma=(0.20,0.30)$
and several correlations (\cref{ch10:ex-two}; compare slides~15--16 of \file{lec13} and the top-right panels of
\file{Smltn_TwoAst}). For $\rho=-1$ the curve reaches the return axis. The lower part of each curve, below the
minimum-variance point, is dominated.}\label{ch10:fig-two}
\end{figure}

\begin{intuition}[Correlation as an angle]
Write $\sigma_w^2=|(1-w)\bm r_1+w\bm r_2|^2$ with vectors $\bm r_1,\bm r_2$ of lengths $\sigma_1,\sigma_2$ and
angle $\theta$, $\cos\theta=\rho$: the volatility of a portfolio is the \emph{length} of a weighted sum of two
vectors. Perfect correlation means parallel vectors (lengths add), $\rho=-1$ antiparallel (lengths cancel), and $\rho=0$
orthogonal (Pythagoras, $\sigma_w^2=(1-w)^2\sigma_1^2+w^2\sigma_2^2$). Diversification is the triangle inequality
being strict. The minimum-variance portfolio is the foot of the perpendicular from the origin to the segment
joining the tips of the two vectors.
\end{intuition}


% ==================================================================
\section{Markowitz mean--variance optimisation}\label{ch10:sec-mvo}

Now let there be $m$ risky assets with return vector $\bm R$, $\E[\bm R]=\bm\mu$, $\Cov\bm R=\Sigma$ (positive
definite) in a \emph{single period}. A portfolio is a vector $\bm w$ of fractions of the wealth invested in each
asset, $\bm w\T\bm 1=1$ (one unit of capital). Its return $R_{\bm w}=\bm w\T\bm R$ has
\begin{equation}\label{ch10:eq-mv}
  \mu_{\bm w}=\bm w\T\bm\mu,\qquad \sigma_{\bm w}^2=\bm w\T\Sigma\bm w
\end{equation}
(\cref{ch03:prop-portvar}) \ts{13}{14}{3:11}\ts{13}{14}{4:16}. Investors prefer higher $\mu_{\bm w}$ and lower
$\sigma^2_{\bm w}$, so portfolios are compared through the pair $(\mu_{\bm w},\sigma_{\bm w}^2)$
\ts{13}{14}{23:14}. Markowitz's idea (1952) is to pose this as a convex optimisation problem.

\begin{definition}[Markowitz problem I]\label{ch10:def-markowitz}
For a target return $\mu_0$, choose $\bm w\in\R^m$ (short positions allowed) to
\[
  \text{minimise } \tfrac12\bm w\T\Sigma\bm w\quad\text{subject to}\quad \bm w\T\bm\mu=\mu_0,\ \ \bm w\T\bm 1=1 .
\]
\end{definition}
The factor $\frac12$ only simplifies the derivatives \ts{13}{14}{24:14}. The objective is strictly convex ($\Sigma\succ0$)
and the constraints are linear, so a stationary point of the Lagrangian is the unique global minimum
\ts{13}{14}{25:20}.

\begin{theorem}[Efficient portfolios without a riskless asset]\label{ch10:thm-mvo}
Put
\[
  a=\bm\mu\T\Sigma^{-1}\bm\mu,\qquad b=\bm\mu\T\Sigma^{-1}\bm 1,\qquad c=\bm 1\T\Sigma^{-1}\bm 1,\qquad
  \Delta=ac-b^2 .
\]
If $\bm\mu$ is not a multiple of $\bm 1$ then $\Delta>0$ and problem~I has the unique solution
\begin{equation}\label{ch10:eq-w0}
  \bm w_0=\lambda_1\Sigma^{-1}\bm\mu+\lambda_2\Sigma^{-1}\bm 1,\qquad
  \begin{pmatrix}\lambda_1\\ \lambda_2\end{pmatrix}=\begin{pmatrix}a&b\\ b&c\end{pmatrix}^{-1}\begin{pmatrix}\mu_0\\ 1\end{pmatrix}
  =\frac1\Delta\begin{pmatrix}c\mu_0-b\\ a-b\mu_0\end{pmatrix},
\end{equation}
whose variance is
\begin{equation}\label{ch10:eq-frontier}
  \sigma_0^2=\bm w_0\T\Sigma\bm w_0=\frac{c\mu_0^2-2b\mu_0+a}{\Delta}
  =\frac1c+\frac{c}{\Delta}\Bigl(\mu_0-\frac bc\Bigr)^2 .
\end{equation}
\end{theorem}
\begin{proof}
The Lagrangian is $L=\frac12\bm w\T\Sigma\bm w+\lambda_1(\mu_0-\bm w\T\bm\mu)+\lambda_2(1-\bm w\T\bm 1)$
\ts{13}{14}{24:14}. Setting $\partial L/\partial\bm w=\Sigma\bm w-\lambda_1\bm\mu-\lambda_2\bm 1=\bm 0$ gives
$\bm w=\lambda_1\Sigma^{-1}\bm\mu+\lambda_2\Sigma^{-1}\bm 1$. Substituting in the two constraints gives the linear
system $\mu_0=a\lambda_1+b\lambda_2$, $1=b\lambda_1+c\lambda_2$, that is \eqref{ch10:eq-w0} \ts{13}{14}{26:26}.
The matrix $\begin{psmallmatrix}a&b\\b&c\end{psmallmatrix}$ is the Gram matrix of $\bm\mu,\bm 1$ for the inner
product $\langle\bm x,\bm y\rangle=\bm x\T\Sigma^{-1}\bm y$; its determinant $\Delta$ is positive by the strict
Cauchy--Schwarz inequality, because $\bm\mu$ and $\bm 1$ are linearly independent. Finally
$\sigma_0^2=\bm w_0\T\Sigma\bm w_0=\lambda_1\bm w_0\T\bm\mu+\lambda_2\bm w_0\T\bm 1=\lambda_1\mu_0+\lambda_2$, which
equals \eqref{ch10:eq-frontier}; completing the square in $\mu_0$ gives the second form.
\end{proof}

Thus $\sigma_0^2$ is a parabola in the target return \ts{13}{14}{27:32}: the smallest variance obtainable
with mean $\mu_0$, and the set of all pairs $(\sigma_0,\mu_0)$ is a hyperbola in the volatility--return plane.

\begin{corollary}[Global minimum-variance portfolio and the frontier]\label{ch10:cor-gmv}
The parabola \eqref{ch10:eq-frontier} is smallest at $\mu_0=b/c$, where
\[
  \bm w_g=\frac{\Sigma^{-1}\bm 1}{\bm 1\T\Sigma^{-1}\bm 1},\qquad \sigma_g^2=\frac1c,\qquad \mu_g=\frac bc .
\]
The \emph{efficient frontier} is the upper branch $\mu_0\ge b/c$ (portfolios below it are dominated by the
portfolio with the same variance and larger return); the whole curve is a hyperbola with asymptotes
$\mu=b/c\pm\sqrt{\Delta/c}\,\sigma$. Along the efficient branch $\mu$ is an increasing concave function of $\sigma$, and no portfolio lies to the left of the hyperbola.
\end{corollary}

\subsection{The three equivalent problems}
The slides pose the same optimisation in three ways (slide~7) \ts{13}{14}{28:36}\ts{13}{14}{29:42}:
\begin{itemize}
  \item \textbf{Problem I}: minimise the variance for a given mean $\mu_0$;
  \item \textbf{Problem II}: maximise $\bm w\T\bm\mu$ for a given variance $\sigma_0^2$, with $\bm w\T\bm 1=1$;
  \item \textbf{Problem III}: for a risk-aversion parameter $\lambda\ge0$ maximise
        $\bm w\T\bm\mu-\frac\lambda2\bm w\T\Sigma\bm w$ with $\bm w\T\bm 1=1$.
\end{itemize}
All three have the same first-order conditions, and the \emph{efficient frontier} is the trace of the solutions
when $\mu_0$ (I), $\sigma_0^2$ (II) or $\lambda$ (III) varies. Problem~III can be solved in one line, and it shows the
structure of the frontier. The condition $\bm\mu-\lambda\Sigma\bm w-\gamma\bm 1=\bm 0$ (with multiplier $\gamma$
for the budget) gives $\bm w=\lambda^{-1}\Sigma^{-1}(\bm\mu-\gamma\bm 1)$, and $\bm 1\T\bm w=1$ fixes $\gamma=(b-\lambda)/c$:
\begin{equation}\label{ch10:eq-twofund}
  \bm w(\lambda)=\bm w_g+\frac1\lambda\,\bm v,\qquad \bm v=\Sigma^{-1}\Bigl(\bm\mu-\frac bc\bm 1\Bigr),\qquad
  \bm 1\T\bm v=0 .
\end{equation}
Every efficient portfolio is the minimum-variance portfolio plus a multiple $1/\lambda$ of one fixed
\emph{zero-cost} (long/short) portfolio $\bm v$. Hence $\mu=b/c+(a-b^2/c)/\lambda$ and
$\sigma^2=1/c+(a-b^2/c)/\lambda^2$, since $\bm v\T\Sigma\bm v=a-b^2/c$ and $\bm w_g\T\Sigma\bm v=0$; eliminating $\lambda$
recovers \eqref{ch10:eq-frontier}. Consequently any two distinct efficient portfolios span the whole frontier (the
\emph{two-fund theorem}, not stated in class). The risk aversion $\lambda$ is the Arrow--Pratt index of
\cref{ch10:sec-utility}: a very risk-averse investor ($\lambda\to\infty$) holds $\bm w_g$.

\begin{example}[Nine sector ETFs, unconstrained (computed from the rounded case-study inputs)]\label{ch10:ex-etf-unc}
Take the annualised means, volatilities and correlations of the nine SPDR sector ETFs for 2009--2013 printed in
Case Study~6 (\cref{ch10:tab-etf}). Then $a=3.44$, $b=11.44$, $c=88.04$, $\Delta=171.9$. The global minimum-variance
portfolio has $\sigma_g=1/\sqrt c=10.7\%$ and $\mu_g=b/c=13.0\%$, better in volatility than every single ETF
(the lowest is consumer staples, $12\%$, but its mean is $15\%$), and the asymptotic slope is
$\sqrt{\Delta/c}=1.40$. But the weights, in the order XLB, XLV, XLP, XLY, XLE, XLF, XLI, XLK, XLU, are
\[
  \bm w_g\approx(0.09,\ 0.14,\ 0.94,\ -0.17,\ -0.15,\ -0.02,\ -0.27,\ 0.26,\ 0.18),
\]
with $94\%$ in consumer staples and a $27\%$ short position in industrials: an unconstrained optimum is not
what an investor would hold. This is the reason for the constraints of \cref{ch10:sec-constr}.
\end{example}

% ==================================================================
\section{A riskless asset: tangency portfolio and capital market line}\label{ch10:sec-rf}

Add a riskless asset $0$ with constant return $R_0\equiv r_0$ (zero variance, uncorrelated with everything)
\ts{13}{14}{31:55}. The investor puts $\bm w$ in the risky assets and $1-\bm w\T\bm 1$ in the riskless one, which
may be negative (borrowing at $r_0$). There is no budget constraint left on $\bm w$, and
\begin{equation}\label{ch10:eq-mvrf}
  \mu_{\bm w}=\bm w\T\bm\mu+(1-\bm w\T\bm 1)r_0=r_0+\bm w\T(\bm\mu-r_0\bm 1),\qquad \sigma_{\bm w}^2=\bm w\T\Sigma\bm w .
\end{equation}
\emph{Geometric picture.} Two assets are enough: combining the riskless asset with a risky portfolio~$P$ gives the
line through $(0,r_0)$ and $(\sigma_P,\mu_P)$; the higher the slope, the better. The best line is the tangent to
the efficient frontier of the risky assets \ts{24}{13}{35:50}\ts{13}{14}{33:01}, which lies above the frontier
because we can now hold a fraction of a risky portfolio in cash \ts{13}{14}{35:05}. The mathematics confirms it.

\begin{theorem}[Optimal portfolios with a riskless asset; Tobin separation]\label{ch10:thm-tobin}
Let $\bm\eta=\bm\mu-r_0\bm 1$ (excess returns) and $d^2=\bm\eta\T\Sigma^{-1}\bm\eta=a-2r_0b+r_0^2c$. The problem
$\min\frac12\bm w\T\Sigma\bm w$ subject to $\mu_{\bm w}=\mu_0$ has the solution
\begin{equation}\label{ch10:eq-rfsol}
  \bm w_0=\lambda_1\Sigma^{-1}\bm\eta,\qquad \lambda_1=\frac{\mu_0-r_0}{d^2},\qquad
  \sigma_0=\frac{|\mu_0-r_0|}{d}.
\end{equation}
Assume $\bm 1\T\Sigma^{-1}\bm\eta=b-r_0c>0$ (i.e.\ $r_0<\mu_g$) and define the \emph{market (tangency) portfolio}
\begin{equation}\label{ch10:eq-market}
  \bm w_M=\frac{\Sigma^{-1}\bm\eta}{\bm 1\T\Sigma^{-1}\bm\eta},\qquad
  \mu_M=\frac{a-r_0b}{b-r_0c},\qquad \sigma_M=\frac{d}{b-r_0c}.
\end{equation}
Then every optimal portfolio with $\mu_0>r_0$ invests in the risky assets in the proportions of $\bm w_M$: it holds
the fraction
\begin{equation}\label{ch10:eq-wm}
  m_{\bm w}=\bm w_0\T\bm 1=\frac{\mu_0-r_0}{\mu_M-r_0}
\end{equation}
of its wealth in $\bm w_M$ and $1-m_{\bm w}$ in the riskless asset \ts{13}{14}{41:30}\ts{13}{14}{42:31}.
Its return and risk are $\mu_P=r_0+m_{\bm w}(\mu_M-r_0)$ and $\sigma_P=m_{\bm w}\sigma_M$.
\end{theorem}
\begin{proof}
The Lagrangian $L=\frac12\bm w\T\Sigma\bm w+\lambda_1[(\mu_0-r_0)-\bm w\T\bm\eta]$ gives
$\Sigma\bm w=\lambda_1\bm\eta$, i.e.\ $\bm w=\lambda_1\Sigma^{-1}\bm\eta$, and the constraint
$\bm w\T\bm\eta=\mu_0-r_0$ gives $\lambda_1d^2=\mu_0-r_0$ \ts{13}{14}{35:05}\ts{13}{14}{36:08}. The variance is
$\lambda_1^2\bm\eta\T\Sigma^{-1}\bm\eta=\lambda_1^2d^2=(\mu_0-r_0)^2/d^2$. Since $\bm w_0=\lambda_1(b-r_0c)\bm w_M$,
the risky weights are proportional to $\bm w_M$ with total weight $m_{\bm w}=\lambda_1(b-r_0c)$; \eqref{ch10:eq-market}
follows from $\mu_M=\bm w_M\T\bm\mu$, $\sigma_M^2=\bm\eta\T\Sigma^{-1}\bm\eta/(\bm 1\T\Sigma^{-1}\bm\eta)^2$, and
$\mu_M-r_0=d^2/(b-r_0c)$ gives \eqref{ch10:eq-wm}.
\end{proof}

\begin{corollary}[Capital market line and price of risk]\label{ch10:cor-cml}
The optimal portfolios lie on the half-line, the \emph{capital market line} (CML)
\begin{equation}\label{ch10:eq-cml}
  \mu_P=r_0+\frac{\mu_M-r_0}{\sigma_M}\,\sigma_P,\qquad \sigma_P\ge0,
\end{equation}
which is tangent to the efficient frontier at $M$. Its slope $(\mu_M-r_0)/\sigma_M=d$ is the \emph{Sharpe ratio} of
the market portfolio, the ``market price of risk'': the compensation per unit of volatility. Investors of all risk
appetites hold the \emph{same} risky portfolio and differ only in how much of it (Tobin, 1958)
\ts{13}{14}{43:42}\ts{13}{14}{44:47}. If borrowing at $r_0$ is possible, $m_{\bm w}>1$ (i.e.\ $\mu_0>\mu_M$) gives points of the CML
beyond $M$, a \emph{leveraged} market portfolio \ts{13}{14}{46:52}.
\end{corollary}

\begin{definition}[Sharpe ratio]
For a portfolio with return $R_P$, the Sharpe ratio is $\mathrm{SR}_P=(\mu_P-r_0)/\sigma_P$
\ts{24}{13}{36:54}\ts{13}{16}{44:21}.
\end{definition}

\begin{proposition}[The tangency portfolio maximises the Sharpe ratio]\label{ch10:prop-sharpe}
For every risky portfolio $\bm w\neq\bm 0$, $(\bm w\T\bm\eta)/\sqrt{\bm w\T\Sigma\bm w}\le d$, with equality iff $\bm w$
is a positive multiple of $\Sigma^{-1}\bm\eta$, i.e.\ of $\bm w_M$. The Sharpe ratio does not change when a portfolio
is mixed with the riskless asset or levered at rate $r_0$.
\end{proposition}
\begin{proof}
Cauchy--Schwarz for the inner product defined by $\Sigma$:
\[(\bm w\T\bm\eta)^2=\bigl(\bm w\T\Sigma\cdot\Sigma^{-1}\bm\eta\bigr)^2
\le(\bm w\T\Sigma\bm w)(\bm\eta\T\Sigma^{-1}\bm\eta).\] Scaling $\bm w\to t\bm w$ multiplies $\mu-r_0$ and $\sigma$ by $t$.
\end{proof}

The last statement is the answer to Xia's question in 2013 (``when you lever up, does the Sharpe ratio change?''):
with $w_1=-1$ in cash and $w_2=2$ in the risky asset the volatility is $2\sigma_2$ and the excess return
$2(R_2-R_1)$, so the ratio $(R_2-R_1)/\sigma_2$ is the same at every point of the line
\ts{13}{16}{56:12}\ts{13}{16}{57:16}\ts{13}{16}{1:01:28}; the 2024 slide~18 draws the same picture,
$\sigma=2\sigma_p$, $R=2R_p-R_f$ \ts{24}{13}{38:56}. Leverage moves the investor along the line, it does not improve it.

\begin{example}[Two assets and a riskless rate]\label{ch10:ex-tobin}
Take $\bm\mu=(0.15,0.25)$, $\Sigma=\mathrm{diag}(0.04,0.09)$ (uncorrelated assets of \cref{ch10:ex-two}) and
$r_0=0.05$. Then $\bm\eta=(0.10,0.20)$, $\Sigma^{-1}\bm\eta=(2.5,\ 2.22)$, $\bm w_M=(0.529,\ 0.471)$,
$\mu_M=19.7\%$, $\sigma_M=17.6\%$ and $d=\sqrt{0.10^2/0.04+0.20^2/0.09}=0.833$; the individual Sharpe ratios are only
$0.50$ and $0.67$. An investor who wants $\mu_0=12\%$ puts $(0.12-0.05)/(0.197-0.05)=47.6\%$ of the wealth in $M$ and
$52.4\%$ in cash; one who wants $\mu_0=20\%$ borrows $2\%$ of the wealth. \sloppy Twice levered ($m_{\bm w}=2$) gives
$\mu=34.4\%$ at $\sigma=35.3\%$, still on the line (\cref{ch10:fig-cml}).
\end{example}

\begin{figure}[ht]
\centering
\begin{tikzpicture}
\begin{axis}[width=0.86\linewidth,height=7cm,xlabel={volatility $\sigma$},ylabel={expected return $\mu$},
  xmin=0,xmax=0.45,ymin=0,ymax=0.40,xtick={0,0.1,0.2,0.3,0.4},scaled ticks=false,/pgf/number format/fixed,grid=major,grid style={black!10},
  tick label style={font=\small},label style={font=\small}]
  \addplot[thick,defcol,dashed] coordinates {(0.3672,0.0900) (0.3354,0.1000) (0.3046,0.1100) (0.2751,0.1200) (0.2474,0.1300) (0.2220,0.1400) (0.2000,0.1500) (0.1825,0.1600) (0.1709,0.1700) (0.1664,0.1808)};
  \addplot[very thick,defcol] coordinates {(0.1664,0.1808) (0.1697,0.1900) (0.1803,0.2000) (0.1970,0.2100) (0.2184,0.2200) (0.2433,0.2300) (0.2707,0.2400) (0.3000,0.2500) (0.3306,0.2600) (0.3622,0.2700) (0.3946,0.2800) (0.4276,0.2900) (0.4610,0.3000)};
  \addplot[thick,thmcol] coordinates {(0,0.05) (0.40,0.3833)} node[pos=0.22,above,sloped,font=\footnotesize]{CML, slope $0.833$};
  \addplot[only marks,mark=*,mark size=2.2pt,thmcol] coordinates {(0.17647,0.19706) (0,0.05) (0.3529,0.3441)};
  \addplot[only marks,mark=square*,mark size=2pt,black] coordinates {(0.1664,0.1808) (0.2,0.15) (0.3,0.25)};
  \node[font=\footnotesize,anchor=east,thmcol] at (axis cs:0.168,0.225) {$M$};
  \node[font=\footnotesize,anchor=south west,thmcol] at (axis cs:0.005,0.055) {$r_0$};
  \node[font=\footnotesize,anchor=north west,thmcol] at (axis cs:0.36,0.335) {leveraged};
  \node[font=\footnotesize,anchor=south west] at (axis cs:0.203,0.153) {asset 1};
  \node[font=\footnotesize,anchor=north west] at (axis cs:0.305,0.245) {asset 2};
  \node[font=\footnotesize,anchor=east] at (axis cs:0.150,0.225) {min.\ variance};\draw[->] (axis cs:0.120,0.215) -- (axis cs:0.1640,0.1830);
    \node[font=\footnotesize,defcol,anchor=east] at (axis cs:0.445,0.19) {frontier (solid: efficient)};
\end{axis}
\end{tikzpicture}
\caption{Frontier of the two assets of \cref{ch10:ex-two} for $\rho=0$ with short sales allowed (upper branch solid,
dominated branch dashed) and the capital market line of \cref{ch10:ex-tobin} with $r_0=5\%$, tangent at $M$.
Points on the CML beyond $M$ are levered.}\label{ch10:fig-cml}
\end{figure}

\begin{intuition}[Why one fund is enough]
The direction $\Sigma^{-1}\bm\eta$ solves the \emph{normal equations} $\Sigma\bm w=\bm\eta$, the same form as
$X\T X\hat{\bm\beta}=X\T\bm y$ of \cref{ch:regression}: it is the portfolio whose covariances with the assets are
proportional to their excess returns, which is exactly the CAPM relation \eqref{ch10:eq-sml} of
\cref{ch10:sec-capm}. Once the best ``direction'' in asset space is fixed, scaling it up or down is only a question of
appetite, like rescaling a normal mode; the riskless asset supplies the one extra degree of freedom (the
length) that the fully invested constraint of \cref{ch10:sec-mvo} did not leave.
\end{intuition}

\begin{coursenote}[Naming and the case $r_0>\mu_g$]
Kempthorne calls $\bm w_M$ ``the market portfolio'' because it is the \emph{fully invested} optimal
portfolio in the sense of $\bm w_M\T\bm 1=1$ \ts{13}{14}{39:15}; the same portfolio is called the ``tangency portfolio''
elsewhere. That it coincides with the cap-weighted market of all assets is a statement of \emph{equilibrium}
(Sharpe 1964, Lintner 1965), proved in \cref{ch03:thm-capm} and used in \cref{ch10:sec-capm}; here it is only the
name of a solution. The slides implicitly assume $b-r_0c>0$; if instead $r_0>\mu_g$ the tangency to the frontier
touches the lower, dominated branch and $\bm w_M$ has a negative Sharpe ratio; the optimal portfolio for $\mu_0>r_0$
is then $-\bm w_M$ scaled (short the ``market''), which is unrealistic and signals an unreasonable $r_0$. The
slides also assume that the borrowing and lending rates coincide; with a higher borrowing rate the CML has a
kink at $M$ and only the segment up to $M$ is tangent to the frontier (not discussed in class).
The key papers listed on slide~16 of \file{lecnote14} are Markowitz (1952), Tobin (1958), Sharpe (1964),
Lintner (1965) and Fama (1970) \ts{13}{14}{47:54}; Kempthorne remarks that virtually all of them have a Nobel prize,
and that Fama had been awarded it two weeks earlier \ts{13}{14}{48:55}.
\end{coursenote}

% ==================================================================
\section{Beta, the security market line and the CAPM}\label{ch10:sec-capm}

The mean--variance solution also contains the CAPM, by a route different from the expected-utility derivation of
\cref{ch03:sec-capm}. This derivation is not given in either lecture, but it explains the slides on
$\alpha,\beta$ (2024 slide~17, 2013 L16).
\begin{proposition}[Security market line from the tangency portfolio]\label{ch10:prop-sml}
Let $M$ be the tangency portfolio \eqref{ch10:eq-market}. For every asset $i$ (and every portfolio),
\begin{equation}\label{ch10:eq-sml}
  \mu_i-r_0=\beta_i(\mu_M-r_0),\qquad \beta_i=\frac{\Cov(R_i,R_M)}{\sigma_M^2}.
\end{equation}
\end{proposition}
\begin{proof}
By definition $\Sigma\bm w_M=\bm\eta/(\bm 1\T\Sigma^{-1}\bm\eta)$, i.e.\ $\Cov(R_i,R_M)=[\Sigma\bm w_M]_i=
\eta_i\,\kappa$ with $\kappa=1/(\bm 1\T\Sigma^{-1}\bm\eta)$. Multiplying by $w_{M,i}$ and summing gives
$\sigma_M^2=\kappa(\mu_M-r_0)$, hence $\eta_i=\Cov(R_i,R_M)/\kappa=\Cov(R_i,R_M)(\mu_M-r_0)/\sigma_M^2$.
\end{proof}
Under the assumptions of Sharpe and Lintner every investor holds $M$, so $M$ is the market portfolio of all risky
assets and \eqref{ch10:eq-sml} is the CAPM (\cref{ch03:thm-capm}). The regression of $R_i-r_0$ on the market
excess return has slope $\beta_i$ and, by \eqref{ch10:eq-sml}, zero intercept: in the language of \cref{def:alphabeta}
\[
  R_P-r_0=\alpha+\beta\,(R_{BM}-r_0)+\varepsilon,\qquad \beta=\frac{\Cov(R_P,R_{BM})}{\sigma_{BM}^2}=\rho\,\frac{\sigma_P}{\sigma_{BM}}
\]
(2024 slide~17 \ts{24}{13}{37:54}), where a nonzero $\alpha$ measures a departure from the security market line, and
$\sigma_P^2=\beta^2\sigma_{BM}^2+\sigma_\varepsilon^2$ splits the risk into a systematic and an idiosyncratic part
(regression, \cref{ch05:capm}). In the two-asset special case of the lectures (riskless asset plus one risky
asset $2$, $\sigma_P=w_2\sigma_2$) the correlation with asset~2 is $1$ and $\beta=\sigma_P/\sigma_2=w_2$
\ts{13}{16}{1:02:34}\ts{13}{16}{1:03:38}: the 60/40 investor asked how much equity beta the portfolio has needs the
same computation with both assets (\cref{ch10:sec-rp}).

\begin{finance}[Alpha, beta and the Sharpe ratio in practice]
Beta is cheap to buy (an index fund); alpha is what a manager is paid for, and the Sharpe ratio $\mathrm{SR}$
measures return per unit of total risk while the \emph{information ratio} $\alpha/\sigma_\varepsilon$ measures
skill per unit of active risk (not discussed in class). A portfolio manager benchmarked to an index is judged on
$\alpha$ and on the tracking error (\cref{ch10:sec-constr}); a fund with a high Sharpe ratio can be levered to
any desired volatility, so Sharpe ratios are the fair basis for comparing funds of different volatilities.
\end{finance}


% ==================================================================
\section{Diversification and rebalancing: the ``free lunch''}\label{ch10:sec-rebal}

\subsection{How much can diversification remove?}
For the equally weighted portfolio of $n$ assets, with $\bar\sigma^2$ the average variance and $\bar c$ the average
pairwise covariance,
\begin{equation}\label{ch10:eq-ew}
  \Var(R_{\rm ew})=\frac1{n^2}\Bigl(\sum_i\sigma_i^2+\sum_{i\ne j}\Sigma_{ij}\Bigr)=\frac{\bar\sigma^2}n+\Bigl(1-\frac1n\Bigr)\bar c .
\end{equation}
The first term (idiosyncratic risk) vanishes like $1/n$, the second does not: what survives is the average
covariance, the \emph{systematic} risk that no diversification removes (\cref{ch07:ex-equi} for the spectral
version). For uncorrelated assets of equal volatility, $\sigma/\sqrt n$; Exercises~\ref{ch10:ex-ps6-4} and
\ref{ch10:ex-ps6-5} ask for the minimum-variance portfolios in the uncorrelated, equicorrelated and general cases.

\subsection{Rebalancing}
Xia's second illustration of diversification in both years is a two-year, two-asset example
\ts{24}{13}{40:00}\ts{13}{16}{52:00}. Asset~1 doubles in year~1 and halves in year~2; asset~2 halves and then doubles.
Alone, each asset ends where it started (return $0$). Start with $50\%$ in each \ts{24}{13}{41:06}.
\begin{itemize}
  \item Year~1: $0.5\cdot(+100\%)+0.5\cdot(-50\%)=+25\%$, so the value is $1.25$, of which
        $1.00$ ($80\%$) is now asset~1 and $0.25$ ($20\%$) asset~2 \ts{24}{13}{42:13}.
  \item \emph{Without rebalancing}, year~2 gives $0.8\cdot(-50\%)+0.2\cdot(+100\%)=-20\%$, and the two-year return is
        $1.25\times0.8-1=0$: no gain from a diversified portfolio of perfectly negatively correlated assets.
  \item \emph{With rebalancing} to $50/50$ at the end of year~1 (sell $0.375$ of asset~1, buy asset~2: $0.625$ each),
        year~2 gives $0.5\cdot(-50\%)+0.5\cdot(+100\%)=+25\%$: two years at $+25\%$, total $1.25^2=1.5625$
        \ts{24}{13}{42:13}\ts{13}{16}{54:05}\ts{13}{16}{55:08}.
\end{itemize}
Rebalanced, the portfolio has \emph{zero} volatility and a steady $25\%$ return, ``a straight line''.

\begin{figure}[ht]
\centering
\begin{tikzpicture}
\begin{axis}[width=0.62\linewidth,height=5.2cm,xlabel={year},ylabel={value of \$1},xmin=-0.1,xmax=2.1,ymin=0.3,ymax=2.2,
  xtick={0,1,2},grid=major,grid style={black!10},legend style={font=\footnotesize,at={(0.5,1.03)},anchor=south,legend columns=2},
  tick label style={font=\small},label style={font=\small}]
  \addplot[thick,defcol,mark=*] coordinates {(0,1) (1,2) (2,1)};
  \addplot[thick,intcol,mark=*] coordinates {(0,1) (1,0.5) (2,1)};
  \addplot[thick,thmcol,mark=square*,dashed] coordinates {(0,1) (1,1.25) (2,1.0)};
  \addplot[very thick,cscol,mark=square*] coordinates {(0,1) (1,1.25) (2,1.5625)};
  \legend{asset 1,asset 2,50/50 no rebalancing,50/50 rebalanced}
\end{axis}
\end{tikzpicture}
\caption{The two-year example of Xia (2024 slide~20, 2013 L16): each asset ends at $1$, the buy-and-hold $50/50$
portfolio also, the rebalanced one at $1.5625$.}\label{ch10:fig-rebal}
\end{figure}

Two objections are discussed in class \ts{24}{13}{43:16}\ts{24}{13}{44:18}. Is this not a bet on mean reversion,
losing money in a trend? Yes: rebalancing sells the winner and buys the loser, which pays when returns alternate
and loses when they trend (the case of \cref{sec:ewport}, where daily rebalancing to equal weights was not obviously
better). The rule of thumb is therefore \emph{consistency of views}: the equal weights expressed the belief that the
two assets have the same distribution and the same risk; unless the view has changed, letting the weights drift
means holding a non-equal portfolio by accident. ``The only free lunch is diversification, but you have to
rebalance to eat it.''

\begin{proposition}[Excess growth of a constant mix; not discussed in class]\label{ch10:prop-growth}
Let asset $i$ follow a geometric Brownian motion with $\dd S_i/S_i=\mu_i\dd t+\sigma_i\dd W_i$ and
$\Cov(\dd W_i,\dd W_j)=\rho_{ij}\dd t$. Asset $i$ alone grows at the logarithmic rate $\mu_i-\frac12\sigma_i^2$.
A portfolio rebalanced \emph{continuously} to the weights $\bm w$ has
\begin{equation}\label{ch10:eq-growth}
  g(\bm w)=\bm w\T\bm\mu-\tfrac12\bm w\T\Sigma\bm w,\qquad
  g(\bm w)-\sum_iw_i\bigl(\mu_i-\tfrac12\sigma_i^2\bigr)=\tfrac12\Bigl(\sum_iw_i\sigma_i^2-\bm w\T\Sigma\bm w\Bigr)\ge0 .
\end{equation}
\end{proposition}
\begin{proof}
The wealth $V$ of the constant mix satisfies $\dd V/V=\bm w\T\bm\mu\dd t+\bm w\T\bm\sigma\dd\bm W$ (its return is
$\bm w\T$ times the vector of returns, \cref{ch:stochcalc}), so by It\^o's formula
$\dd\log V=(\bm w\T\bm\mu-\tfrac12\bm w\T\Sigma\bm w)\dd t+\dots$. The inequality is $\Var(\sum_iw_iR_i)\le\sum_iw_i\Var R_i$,
i.e.\ convexity of the variance.
\end{proof}
The gain is the ``volatility harvesting'' or ``diversification return''. Two uncorrelated assets with $\sigma=20\%$
and the same $\mu$, kept at $50/50$, grow at $\mu-1\%$ per year, while each alone grows at $\mu-2\%$: the
rebalanced portfolio gains $\sigma^2/4=1\%$ per year over the average of its components. The example above is the
extreme case: $\rho=-1$, growth rates of $0\%$ for each asset, $25\%$ for the mix.

% ==================================================================
\section{Leverage, the 60/40 portfolio and risk parity}\label{ch10:sec-rp}

Two-asset portfolios are important in practice, because managers are often benchmarked to a mix of equities and
bonds, the most famous being the \emph{60/40} portfolio: $60\%$ in equities and $40\%$ in bonds. Any fund is still
compared with it \ts{13}{16}{41:28}\ts{13}{16}{42:34}\ts{13}{16}{47:42}.

\subsection{Weights are not risks}
Equities are much more volatile than bonds, so a $60/40$ portfolio is really an equity portfolio in risk terms.
The idea of \emph{risk parity} (the name was coined in 2005 by Edward Qian, as Xia recalls) is to allocate
\emph{risk} rather than capital: equal risk from each component, instead of weights by market value
\ts{13}{16}{49:54}\ts{13}{16}{50:55}\ts{24}{13}{38:56}. Before 2000, Xia recalls, a passive $60/40$ was the norm; the
2000 equity peak and bond rally and the 2008 crash, when government bonds rallied while equities fell, made
investors question it \ts{13}{16}{48:48}\ts{13}{16}{49:54}. To define the risk contribution of an asset one
needs a decomposition of $\sigma_P=\sqrt{\bm w\T\Sigma\bm w}$; it comes from Euler's theorem, since $\sigma_P$ is
homogeneous of degree one in $\bm w$.

\begin{definition}[Risk contributions]\label{ch10:def-rc}
The \emph{risk contribution} of asset $i$ to a portfolio $\bm w$ is
\[
  \mathrm{RC}_i=w_i\frac{\partial\sigma_P}{\partial w_i}=\frac{w_i(\Sigma\bm w)_i}{\sigma_P}=\frac{w_i\Cov(R_i,R_P)}{\sigma_P},
  \qquad \sum_i\mathrm{RC}_i=\sigma_P .
\]
A portfolio has \emph{equal risk contributions} (risk parity, ERC) if $\mathrm{RC}_i=\mathrm{RC}_j$ for all $i,j$, that is
$w_i(\Sigma\bm w)_i=$ const.
\end{definition}
The identity $\sum_i\mathrm{RC}_i=\sigma_P$ is Euler's theorem, or simply $\sum_iw_i(\Sigma\bm w)_i=\bm w\T\Sigma\bm w=\sigma_P^2$.

\begin{proposition}[Risk parity for two assets and for equal correlations]\label{ch10:prop-rp}
\leavevmode
\begin{enumerate}[(i)]
  \item For two assets with volatilities $\sigma_1,\sigma_2$ and \emph{any} correlation $\rho>-1$ the risk-parity
      weights are $w_1:w_2=1/\sigma_1:1/\sigma_2$ (inverse volatility).
  \item The same holds for $m$ assets with equal
      pairwise correlations, whatever the volatilities.
  \item If, in addition, all assets have the \emph{same Sharpe
      ratio}, this risk-parity portfolio is the tangency portfolio $\bm w_M$ of \cref{ch10:thm-tobin}.
\end{enumerate}
\end{proposition}
\begin{proof}
(i) If $w_1\sigma_1=w_2\sigma_2=:\kappa$ then $w_1(\Sigma\bm w)_1=w_1(w_1\sigma_1^2+w_2\rho\sigma_1\sigma_2)=\kappa^2(1+\rho)$ and
likewise for asset~2.

(ii) With $\Sigma=D C D$, $D=\diag(\sigma_i)$, $C=(1-\rho)I+\rho\bm 1\bm 1\T$, and
$w_i=\kappa/\sigma_i$, i.e.\ $D\bm w=\kappa\bm 1$: $w_i(\Sigma\bm w)_i=w_i\sigma_i(C D\bm w)_i=\kappa^2(C\bm 1)_i=\kappa^2(1+(m-1)\rho)$,
independent of $i$.

(iii) $\mu_i-r_0=s\sigma_i$ gives $\bm\eta=sD\bm 1$ and $\Sigma^{-1}\bm\eta=sD^{-1}C^{-1}\bm 1=
s\,D^{-1}\bm 1/(1+(m-1)\rho)$, proportional to $1/\sigma_i$.
\end{proof}

\begin{example}[60/40 versus risk parity, hypothetical numbers]\label{ch10:ex-rp}
Let equities have volatility $15\%$ and bonds $5\%$, uncorrelated, and let both have a Sharpe ratio of $0.3$
(excess returns $4.5\%$ and $1.5\%$). These numbers are \emph{not} from the course; they are chosen so that the
volatility ratio is $3$, which reproduces the $25/75$ equity/bond risk-parity split quoted in class
\ts{13}{16}{59:21}. Then:
\begin{center}
\begin{tabular}{@{}lrrrrr@{}}
\toprule
portfolio (eq./bond) & $\sigma_P$ & excess ret. & Sharpe & eq.\ risk & bond risk\\
\midrule
$60/40$ & $9.2\%$ & $3.3\%$ & $0.358$ & $95.3\%$ & $4.7\%$\\
risk parity $25/75$ & $5.3\%$ & $2.25\%$ & $0.424$ & $50\%$ & $50\%$\\
risk parity levered $1.74\times$ & $9.2\%$ & $3.9\%$ & $0.424$ & $50\%$ & $50\%$\\
\bottomrule
\end{tabular}
\end{center}
The $60/40$ portfolio takes $95\%$ of its risk from equities. Risk parity has a higher Sharpe ratio (here it is
the tangency portfolio by \cref{ch10:prop-rp}(iii), $0.3\sqrt2=0.424$) but a lower return; to get it back, one
\emph{levers} it: borrowing $74\%$ of the capital at the riskless rate brings the volatility up to that of the $60/40$
portfolio and the excess return to $3.9\%$ against $3.3\%$
\ts{13}{16}{56:12}\ts{13}{16}{59:21}\ts{24}{13}{40:00}. Values computed with the formulas above.
\end{example}
That is the ``levered risk parity'' of 2024 slide~19: a levered bond allocation to achieve an equal risk
contribution from equities and bonds, so that the leverage is what buys back the return. Because Sharpe ratios
are invariant to leverage (\cref{ch10:prop-sharpe}), the gain comes from choosing the point of maximal Sharpe ratio.

\begin{finance}[Limits of risk parity]
Risk parity uses \emph{no return forecasts}, only volatilities and correlations; this is attractive because,
a remark from the audience (apparently by Kempthorne, whom the transcript does not name) put it, returns cannot be
estimated accurately even with more data, whereas volatilities can (\cref{ch10:sec-est}) \ts{13}{16}{1:25:00}. But it does depend on a forecast of volatility from history: asked why risk-parity
funds lost heavily in May--June 2013, Xia explained that they overweight low-volatility bonds, which sold off when
the Fed chairman announced the tapering of quantitative easing \ts{13}{16}{1:14:26}. He then asked whether this
episode proves the approach wrong, or the 2008 crisis proves it right, and concluded that the evidence is
inconclusive: one is ``driving while looking in the rear-view mirror''
\ts{13}{16}{1:15:30}. Leverage carries liquidity risk, not treated in the lecture \ts{13}{16}{59:21}.
\end{finance}

\begin{intuition}[Leverage and beta]
Levering a portfolio by a factor $\ell$ multiplies its beta to every factor, and its volatility, by $\ell$, and
the excess return by $\ell$: the Sharpe ratio is unchanged, like the direction of a vector when it is rescaled.
For $\sigma_P=w_2\sigma_2$ in the two-asset special case, $\beta=\sigma_P/\sigma_2$
\ts{13}{16}{1:03:38}: ``a higher volatility of the portfolio means a higher beta to the risky asset''.
\end{intuition}

% ==================================================================
\section{Utility theory: why mean and variance?}\label{ch10:sec-utility}

Markowitz reduces a portfolio to its mean and variance. When is that legitimate? Kempthorne's answer is the
\emph{von Neumann--Morgenstern theory} of decisions under uncertainty: a rational agent has a utility function
$u$ of wealth and chooses the portfolio $P$ maximising the \emph{expected utility} $\E[u(W)]$ of the terminal
wealth $W=W_0(1+R_P)$ \ts{13}{14}{48:55}\ts{13}{14}{50:00}\ts{13}{14}{51:01}. Sensible preferences give
$u'>0$ (more is better) and $u''<0$ (decreasing marginal utility, risk aversion) \ts{13}{14}{52:01}. The
Arrow--Pratt coefficients are
\begin{align*}
  \lambda_A(W)&=-\frac{u''(W)}{u'(W)}\quad\text{(absolute risk aversion)},\\
  \lambda_R(W)&=-\frac{W\,u''(W)}{u'(W)}\quad\text{(relative)}.
\end{align*}
(Slide~19 of \file{lecnote14}; the same $\lambda_A$ appears in \cref{ch03:sec-capm}, where the CARA case is treated.)

\begin{table}[ht]
\centering
\begin{tabular}{@{}llccl@{}}
\toprule
& $u(W)$ & $\lambda_A(W)$ & $\lambda_R(W)$ & remark\\
\midrule
linear & $a+bW$, $b>0$ & $0$ & $0$ & risk neutral\\
quadratic & $W-\frac12\lambda W^2$, $W<1/\lambda$ & $\dfrac\lambda{1-\lambda W}$ & $\dfrac{\lambda W}{1-\lambda W}$ & mean--variance is exact\\
exponential & $1-e^{-\lambda W}$ & $\lambda$ & $\lambda W$ & CARA\\
power & $W^{1-\lambda}$, $0<\lambda<1$ & $\lambda/W$ & $\lambda$ & CRRA\\
logarithmic & $\ln W$ & $1/W$ & $1$ & CRRA, Kelly\\
\bottomrule
\end{tabular}
\caption{The utility functions of slide~20 of \file{lecnote14} \ts{13}{14}{54:08}, with their risk-aversion coefficients (computed here).
Quadratic utility has \emph{increasing} absolute risk aversion and is not increasing beyond $W=1/\lambda$.}\label{ch10:tab-utility}
\end{table}

\subsection{From expected utility to mean--variance}
\begin{proposition}[Three routes to mean--variance analysis]\label{ch10:prop-mv-utility}
\leavevmode
\begin{enumerate}[(i)]
  \item For any smooth $u$, expanding around $w_*=\E[W]$ gives to second order
      $\E[u(W)]\approx u(w_*)+\tfrac12u''(w_*)\Var W$, i.e.\ up to a positive factor and a constant,
      $\E[W]-\frac12\lambda_A(w_*)\Var W$ \ts{13}{14}{53:06}\ts{13}{14}{54:08}.
  \item For quadratic utility the expected utility is exactly $\mu_W-\frac\lambda2(\mu_W^2+\sigma_W^2)$, a function of the mean
      and variance of $W$ alone \ts{13}{14}{55:09}.
  \item If $W\sim N(\mu_W,\sigma_W^2)$, for \emph{any} increasing concave $u$, $\E[u(W)]$ is increasing in $\mu_W$ and
      decreasing in $\sigma_W^2$, so that only the pair $(\mu_W,\sigma_W)$ matters and the efficient frontier is the set of
      candidates for the optimum; for exponential utility, $\E[u(W)]=1-\exp\{-\lambda(\mu_W-\frac\lambda2\sigma_W^2)\}$.
\end{enumerate}
\end{proposition}
\begin{proof}
(i) Taylor's formula at $w_*$, with $\E[W-w_*]=0$: the first-order term drops, and $u''=-\lambda_Au'$.

(ii) Compute
$\E[W-\frac\lambda2W^2]$.

(iii) Write $\E[u(W)]=\int u(x)\varphi(x;\mu_W,\sigma_W^2)\dd x$; the Gaussian density satisfies the
heat equation $\partial_{\sigma^2}\varphi=\frac12\partial_x^2\varphi$ and $\partial_\mu\varphi=-\partial_x\varphi$, so
integrating by parts, $\partial_{\sigma^2}\E[u]=\frac12\E[u''(W)]<0$ and $\partial_\mu\E[u]=\E[u'(W)]>0$. For the exponential,
use the moment generating function of the normal (\cref{ch03:eq-normalmgf}).
\end{proof}
\begin{intuition}[Variance as the heat-equation time]
For a Gaussian, the variance plays the role of the time in the diffusion equation: $\E[u(W)]$ is $u$ smoothed by
a Gaussian of width $\sigma_W$, and smoothing a concave function can only lower its value.
\end{intuition}
So mean--variance analysis is the right thing to do if the utility is quadratic, \emph{or} if returns are Gaussian
(Kempthorne, 2013: ``the imagination that returns are Gaussian''), at least to first orders otherwise
\ts{13}{14}{55:09}\ts{13}{14}{56:11}. With exponential utility and normal returns, wealth
$W=W_0(1+\bm w\T\bm R)$ has certainty equivalent $W_0\bigl(1+\bm w\T\bm\mu-\frac\lambda2\bm w\T\Sigma\bm w\bigr)$ with
$\lambda=\lambda_AW_0$, exactly \emph{problem III} of \cref{ch10:sec-mvo}; with a riskless asset the solution is
$\bm w=\lambda^{-1}\Sigma^{-1}\bm\eta$, i.e.\ the fraction $m_{\bm w}=(\mu_M-r_0)/(\lambda\sigma_M^2)$ of the wealth in the
market portfolio (from \cref{ch10:thm-tobin}): the risk aversion fixes only \emph{how much}, in accord with
Tobin's separation.

\subsection{Higher moments}
When returns are not Gaussian one goes on in the Taylor series (slides 32--33 of \file{lecnote14})
\ts{13}{14}{56:11}:
\[
  \E[u(W)]=u(w_*)+\tfrac12u''(w_*)\Var W+\tfrac1{3!}u'''(w_*)\,m_3+\tfrac1{4!}u^{(4)}(w_*)\,m_4+\cdots,\qquad m_k=\E[(W-w_*)^k],
\]
with $m_3$ and $m_4$ the third and fourth central moments (the slides call them ``Skew'' and ``Kurtosis''). For
the utilities of \cref{ch10:tab-utility} (exponential, power, log) $u'''>0$ and $u^{(4)}<0$: investors like positive
skewness and dislike large fourth moments. The corresponding optimisation
$\max\ \E[R_P]-\lambda_1\Var R_P+\lambda_2m_3-\lambda_3m_4$ subject to $\bm 1\T\bm w=1$ has no closed form; the slides mention
multi-objective methods and \emph{polynomial goal programming}, and warn that skewness and kurtosis are hard
to estimate (they are sensitive to outliers and need large samples).

% ==================================================================
\section{Sizing: the Kelly criterion and the gain--loss ratio}\label{ch10:sec-kelly}

\subsection{Kelly}
The \emph{Kelly formula} is mentioned by Xia and, at the end of the lecture, by Kempthorne \ts{13}{16}{45:21}\ts{13}{16}{1:26:07}. Consider repeated
independent bets in which a stake gains $g$ per unit staked with probability $p$ and loses $l$ per unit with
probability $q=1-p$. Betting a fraction $f$ of the wealth each time multiplies the wealth by $1+fg$ or $1-fl$, and by the
law of large numbers (\cref{ch03:sec-limits}) $\frac1n\log W_n\to G(f)=p\log(1+fg)+q\log(1-fl)$ almost surely. The
fraction that maximises the long-run growth rate is
\begin{equation}\label{ch10:eq-kelly}
  f^*=\frac pl-\frac qg=\frac{pg-ql}{gl},
\end{equation}
since $G'(f)=pg/(1+fg)-ql/(1-fl)=0$ gives $pg(1-fl)=ql(1+fg)$, i.e.\ $f\,gl(p+q)=pg-ql$. For an even-money bet
($g=l=1$), $f^*=p-q=2p-1$: with probability $\frac12$ one bets nothing, with $p=1$ everything, and with $p=0$ one bets
everything \emph{on the other side}, as Xia says \ts{13}{16}{45:21}\ts{13}{16}{46:23}. It is expected-log utility
(\cref{ch10:tab-utility}), the ``Kelly criterion'' of Kelly, Shannon and Thorp, criticised by Samuelson in the
debate recalled in the case of \cref{ch03:sec-capm} and in \emph{Fortune's Formula}; the co-lecturer, speaking from the audience,
stresses that, unlike the single-period theory of this chapter, it deals with \emph{multi-period} investment
\ts{13}{16}{1:26:07}\ts{13}{16}{1:27:08}. For one asset with Gaussian excess return $\mu$ and variance
$\sigma^2$ per period the growth rate of a fraction $f$ is $r_0+f\mu-\frac12f^2\sigma^2$
(\cref{ch10:prop-growth}), maximised at $f^*=\mu/\sigma^2$, the Sharpe ratio divided by the volatility
(not discussed in class); Kelly is problem III with $\lambda=1$.

\subsection{Xia's gain--loss ratio}
Xia argues that volatility is a poor risk measure because it is symmetric; a long out-of-the-money call
\emph{wants} high volatility, a short option position wants low volatility \ts{24}{13}{47:28}, and the Sharpe ratio
inherits the flaw; the Sortino ratio separates the downside but does not tell how much to invest
\ts{24}{13}{48:30}. He proposes to summarise the distribution of the profit and loss $R$ of an investment
(per period, or of the terminal outcome) by the \emph{expected gain} and \emph{expected loss}, both positive numbers,
\[
  G=\E[R^+],\qquad L=\E[R^-],\qquad \text{so}\qquad G-L=\E[R],\quad G+L=\E[|R|] ,
\]
and to rank investments by the ratio
\begin{equation}\label{ch10:eq-gl}
  \frac{G-L}{G+L}=1-\frac{2L}{G+L}=\frac{\E[R]}{\E[|R|]}\in[-1,1],
\end{equation}
which he interprets directly as a position size (2024 slide~24) \ts{24}{13}{48:30}\ts{24}{13}{49:33}. In the
$(L,G)$ plane every investment is a point; the $45^\circ$ line is $G=L$ (zero expected return) and the
ratio is the slope of the ray from the origin: one wants points far above the diagonal
(slide~24: ``$G\gg L$'') and maximises $G$ for given $L$, which replaces the volatility budget by a
\emph{loss} budget \ts{24}{13}{50:37}\ts{24}{13}{51:37}. For a discrete payoff, $G$ and $L$ are sums of probability
times payoff over winning and losing outcomes; in the continuous case they are integrals of the payoff against the
density over $R>0$ and $R<0$, or against a target return other than $0$ \ts{24}{13}{50:37}. Estimation needs only
``half distributions'', no more work than a mean and a standard deviation \ts{24}{13}{52:41}\ts{24}{13}{53:47};
the gain and loss of a combination of two investments depends on how they move together and is found by
Monte Carlo simulation of the joint distribution \ts{24}{13}{54:54}. The \emph{2024 Investment Game} of \cref{ex:game}
is the empirical version: for the chosen stock, sum the daily gains and losses over two months and compute
$(G-L)/(G+L)$; picking a very volatile stock raises $G-L$ but also the denominator, so the ratio ranks
investments of different volatilities on a common scale \ts{24}{13}{51:37}.

\begin{proposition}[Properties of the ratio \eqref{ch10:eq-gl}; not proved in class]\label{ch10:prop-gl}
\leavevmode
\begin{enumerate}[(i)]
  \item For an even-money bet ($R=\pm1$ with probabilities $p,q$) the ratio is $p-q$, the Kelly fraction \eqref{ch10:eq-kelly};
      for $g\neq l$ the two differ (for $g=2$, $l=1$, $p=\frac12$: Kelly $f^*=0.25$, ratio $=1/3$). So ``the same as Kelly
      for binary bets'' (slide~24) holds for even-money bets, or as an approximation.
  \item If $R\sim N(m,s^2)$ the ratio equals $S/h(S)$ with $S=m/s$ and $h(S)=2\varphi(S)+S(2\Phi(S)-1)$, an increasing function
      of the Sharpe ratio: for symmetric returns it gives the same ranking as the Sharpe ratio.
  \item For skewed payoffs the two rankings differ: $R\equiv2$ w.p.\ $\frac12$ and $0$ w.p.\ $\frac12$ has mean $1$ and standard deviation $1$,
      so does $N(1,1)$, and the ratios are $1$ (never a loss) and $0.857$.
\end{enumerate}
\end{proposition}
\begin{proof}
(i) $G=p$, $L=q$, $G+L=1$. For the Kelly value use \eqref{ch10:eq-kelly}.

(ii) For $R=m+sZ$,
$\E[|R|]=s\bigl[2\varphi(S)+S(2\Phi(S)-1)\bigr]$ (fold the normal), and $h'(S)=2\Phi(S)-1$, so
$(S/h)'=(h-Sh')/h^2=2\varphi(S)/h^2>0$.

(iii) $G=1$, $L=0$; for $N(1,1)$, $h(1)=1.1666$ and $1/h(1)=0.857$.
\end{proof}

% ==================================================================


\section{Portfolio constraints}\label{ch10:sec-constr}

Real portfolios are subject to constraints, which Kempthorne lists (slides 22--25 of \file{lecnote14})
\ts{13}{14}{57:16}\ts{13}{14}{58:20}. With $\bm w$ the target portfolio, $\bm w_0$ the current one, $\bm x=\bm w-\bm w_0$ the
trade and $\bm w_B$ a benchmark portfolio:
\begin{itemize}
  \item \textbf{Long-only}: $w_j\ge0$. \textbf{Holding constraints}: $L_j\le w_j\le U_j$, for instance a cap
        on each position; in a medium-frequency equity strategy, a position must not exceed a day's trading
        volume, or it could not be liquidated \ts{13}{14}{58:20}.
  \item \textbf{Turnover}: $|x_j|\le U_j$ for each asset and $\sum_j|x_j|\le U_*$ overall, because one cannot
        trade arbitrarily much in one period.
  \item \textbf{Benchmark exposure}: $\sum_i|w_i-w_{B,i}|\le U_B$, to remain close to an index (S\&P~500,
        NASDAQ-100, Russell) while trying to beat it \ts{13}{14}{59:30}.
  \item \textbf{Tracking error}: $\mathrm{TE}_P=R_P-R_B=(\bm w-\bm w_B)\T\bm R$, so
        $\Var(\mathrm{TE}_P)=(\bm w-\bm w_B)\T\Sigma(\bm w-\bm w_B)\le\bar\sigma^2_{\mathrm{TE}}$
        \ts{13}{14}{1:00:31}.
  \item \textbf{Factor constraints}: for a factor model $R_{i,t}=\alpha_i+\sum_k\beta_{i,k}f_{k,t}+\epsilon_{i,t}$
        (\cref{ch07:sec-lfm}), the exposure to factor $k$ is $\sum_i\beta_{i,k}w_i$ and one imposes
        $|\sum_i\beta_{i,k}w_i|\le U_k$; many strategies \emph{neutralise} all exposures, $\sum_i\beta_{i,k}w_i=0$
        \ts{13}{14}{1:01:33}\ts{13}{14}{1:02:33}.
  \item \textbf{Minimum transaction or holding size} and \textbf{integer constraints}: equities trade in round lots
        and one share of a stock worth several hundred dollars is a significant amount compared with one of a
        \$50 stock, which matters for small portfolios but not for large ones \ts{13}{14}{1:02:33}\ts{13}{14}{1:03:33}.
\end{itemize}
Constraints of the first four types are linear or quadratic in $\bm w$: slide~25 writes general linear constraints
$A_w\bm w\le\bm u_w$, $A_x\bm x\le\bm u_x$, $A_B(\bm w-\bm w_B)\le\bm u_B$ and quadratic ones $\bm w\T Q_w\bm w\le q_w$,
$\bm x\T Q_x\bm x\le q_x$, $(\bm w-\bm w_B)\T Q_B(\bm w-\bm w_B)\le q_B$. The optimisation is then the same as before with
additional Lagrange multipliers, that is a \emph{convex quadratic programme} \ts{13}{14}{1:03:33}. Integer constraints make
it non-convex.

\begin{proposition}[Structure of the constrained solution; not discussed in class]\label{ch10:prop-constr}
Consider the problem I with additional linear inequality constraints.
\begin{enumerate}[(i)]
  \item The optimal value $\sigma_0^2(\mu_0)$ is
      never smaller than the unconstrained one \eqref{ch10:eq-frontier}, and it increases as constraints are added or tightened.
  \item The Karush--Kuhn--Tucker conditions with a fixed set of active constraints are \emph{linear} in
      $(\bm w,\text{multipliers})$ and in $\mu_0$; therefore the optimal weights are piecewise affine in the target return
      $\mu_0$, and the pieces end when a new constraint becomes active or a multiplier changes sign.
\end{enumerate}
\end{proposition}
\begin{proof}
(i) A smaller feasible set has a larger minimum.

(ii) With the active constraints imposed as equalities,
the first-order conditions are the linear system of \cref{ch10:thm-mvo} enlarged by the active rows, whose
solution is affine in the right-hand side $(\mu_0,1,\text{bounds})$; strict convexity makes it unique.
\end{proof}
That is exactly what the plots of the case study show (\cref{ch10:sec-case}): weights that grow along straight
lines in the target return, kinks when a weight reaches its cap, and new assets entering. Below the first
kink (only the constraints $w_j\ge0$ active) the solution is a \emph{scaled} fixed portfolio: ``they enter in the same
proportion'', the de-levered tangency portfolio of \cref{ch10:thm-tobin} computed on the assets that enter, i.e.\ the
\emph{long-only} tangency portfolio, not the unconstrained one, which has short positions (Case Study~6, and Kempthorne
\ts{13}{14}{1:09:02}).

\begin{finance}[Benchmarked management]
A manager measured against an index does not solve problem~I but the same quadratic problem in the \emph{active
weights} $\bm x_B=\bm w-\bm w_B$: maximise $\bm x_B\T\bm\mu-\frac\lambda2\bm x_B\T\Sigma\bm x_B$, the tracking-error
variance being the risk. The efficient frontier is then drawn in the plane (tracking error, active return), and its
slope is the information ratio. Factor-neutrality constraints remove market, sector or style bets so that the
return is pure stock selection.
\end{finance}

% ==================================================================
\section{Estimation error and the fragility of the optimiser}\label{ch10:sec-est}

Everything above takes $\bm\mu$ and $\Sigma$ as known, ``the problem where we assume basically certainty''
\ts{13}{14}{2:07}. In practice they are estimated, and the estimates are noisy \ts{13}{14}{20:08}. Xia's first
objection to modern portfolio theory is exactly that: the capital market assumptions of return,
volatility and correlation are not reliable, and the mean--variance solution is very sensitive to them, is
``unstable'' and has ``unlimited'' solutions unless bounded by artificial constraints
\ts{24}{13}{45:22}\ts{24}{13}{46:27}.

\subsection{How noisy is a sample mean?}
Suppose returns are i.i.d.\ Gaussian with annual mean $\mu$ and volatility $\sigma$ and are observed for $T$ years, at any
frequency. The sample mean has standard error $\sigma/\sqrt T$ \emph{independent of the sampling frequency}, since only
the total time span matters, whereas the sample volatility, computed from $n$ increments, has relative standard error
$1/\sqrt{2n}$ (\cref{ch03:sec-limits}, \cref{ch:volatility}). For 500 weekly returns ($T=9.6$ years, $n=500$):
\begin{center}
\begin{tabular}{@{}lccc@{}}
\toprule
& true value & standard error & sample values in \file{Smltn_TwoAst}\\
\midrule
mean of asset 1 & $15\%$ & $\sigma_1/\sqrt T=6.4\%$ & $0.085$--$0.21$\\
mean of asset 2 & $25\%$ & $\sigma_2/\sqrt T=9.7\%$ & $0.177$--$0.343$\\
vol.\ of asset 1 & $20\%$ & $3.2\%$ of itself ($0.6\%$) & $0.191$--$0.200$\\
vol.\ of asset 2 & $30\%$ & $3.2\%$ of itself ($0.9\%$) & $0.287$--$0.291$\\
\bottomrule
\end{tabular}
\end{center}
This is what Kempthorne shows by simulation: sample means may be far from the population values while volatilities
and correlations are estimated much better \ts{13}{14}{21:08}\ts{13}{14}{22:11}, and what the audience remark of
2013 L16 says about returns versus volatilities \ts{13}{16}{1:25:00}. In finance ``more data'' means longer periods
(the mean does not improve by sampling faster), while regime changes limit how far back one can go: the choice of
the estimation period is itself an issue \ts{13}{14}{1:20:06}. With the five-year sector-ETF sample of
Case Study~6 the mean of a sector with $\sigma=20\%$ has a standard error of $9\%$, larger than most differences
between the sector means ($10\%$ to $24\%$).

\subsection{What the optimiser does with the errors}
Because $\bm w\propto\Sigma^{-1}\bm\eta$ and highly correlated assets make $\Sigma$ nearly singular, errors in
the means are amplified in the weights. On the rounded nine-ETF inputs of \cref{ch10:tab-etf} (computed here), the
tangency portfolio for $r_0=0$ is extremely leveraged, and \emph{very small changes of the means, of one
percentage point per asset}, change it a lot:
\begin{center}
\footnotesize
\begin{tabular}{@{}lrrrrrrrrrr@{}}
\toprule
weights & XLB & XLV & XLP & XLY & XLE & XLF & XLI & XLK & XLU & $\sum|\Delta w|$\\
\midrule
base means & $-0.32$ & $0.59$ & $0.58$ & $1.49$ & $-0.06$ & $-0.36$ & $-0.45$ & $-0.10$ & $-0.35$ & \\
XLY mean $+1$ pp & $-0.34$ & $0.63$ & $0.53$ & $1.73$ & $-0.05$ & $-0.39$ & $-0.54$ & $-0.18$ & $-0.38$ & $0.58$\\
$\pm1$ pp draw A & $-0.32$ & $0.75$ & $0.28$ & $1.19$ & $-0.09$ & $-0.35$ & $-0.29$ & $0.07$ & $-0.24$ & $1.24$\\
$\pm1$ pp draw B & $-0.57$ & $0.68$ & $0.42$ & $2.07$ & $-0.05$ & $-0.55$ & $-0.19$ & $-0.33$ & $-0.48$ & $1.91$\\
\bottomrule
\end{tabular}
\end{center}
(Draw A is $+,+,-,-,+,+,+,+,+$ pp and draw B $-,-,-,+,-,-,+,-,-$ pp for XLB to XLU; the weights sum to $1$.) For comparison,
the sum of absolute weights is $4.3$: a $\pm1$pp error in each mean, well inside the statistical
error, moves $29\%$ and $44\%$ of that gross exposure (draws A and B). In this respect the caps of \cref{ch10:sec-constr} act as
regularisation, ``artificial constraints'' that bound the solutions \ts{24}{13}{45:22}. The minimum-variance portfolio
$\bm w_g$ does not use $\bm\mu$ at all, which is why it is popular.

\subsection{Remedies}
The slides list the standard alternatives \ts{13}{14}{1:20:06} (slide~27 of \file{lecnote14}): the sample means and
covariance are the least-squares, unbiased and, under Gaussian assumptions, maximum likelihood estimates, but
one may instead use \emph{exponentially weighted moving averages} (\cref{ch:volatility}), \emph{dynamic factor
models}, and optimisation with a simple model for $\Sigma$, notably Sharpe's \emph{single-index model}
$R_i=\alpha_i+\beta_iR_m+\epsilon_i$ (with uncorrelated $\epsilon_i$), which gives $\Sigma=\sigma_m^2\bm\beta\bm\beta\T+
\diag(\sigma^2_{\epsilon,i})$, or the \emph{constant correlation model}, $\Sigma_{ij}=\bar\rho\sigma_i\sigma_j$. A factor model
estimates $2m+1$ parameters instead of $m(m+1)/2$ (for $m=500$ stocks, $1001$ instead of $125{,}250$), and ``results in
more precise inputs to the optimisation'' \ts{13}{14}{1:20:06}; see \cref{ch07:sec-lfm,ch:regularized} (shrinkage,
regularised risk models) and \cref{ch05:sec-shrink}. The returns are the weak point: forecasting them is the actual
job of an active manager, as an audience member observed in 2013 \ts{13}{16}{1:22:59}\ts{13}{16}{1:24:00}.

\begin{finance}[Why the frontier is estimated with errors: ``error maximisation'']
An optimiser puts most weight on the assets whose estimated means are too high and whose estimated variances are too
low, which is precisely where estimation errors are largest: the out-of-sample performance of the
mean--variance solution is typically worse than the in-sample frontier suggests (not discussed in class).
\end{finance}

% ==================================================================
\section{Other risk measures}\label{ch10:sec-risk}

Volatility is only one measure of risk, and Kempthorne closes 2013 L14 by showing how the machinery extends
\ts{13}{14}{1:20:06}\ts{13}{14}{1:21:13} (slides 29--33). Let $R_P=\bm w\T\bm R$.
\begin{itemize}
  \item \textbf{Mean absolute deviation}, $\mathrm{MAD}=\E[|\bm w\T(\bm R-\bm\mu)|]$, leads to a linear programme with linear or
        quadratic constraints. Where the variance is the natural scale for Gaussian returns, MAD is natural for
        heavier-tailed, e.g.\ exponential (Laplace), returns \ts{13}{14}{1:21:13}.
  \item \textbf{Semi-variance}, $\E[\min(R_P-\E[R_P],0)^2]$, the downside variance, introduced by Markowitz himself:
        it controls downside rather than upside risk \ts{13}{14}{1:21:13}. Its square root, relative to a target,
        is the denominator of the Sortino ratio \ts{24}{13}{48:30}.
  \item \textbf{Value at risk}, $\mathrm{VaR}_{1-\epsilon}(R_P)=\min\{r:\Prob(R_P\le-r)\le\epsilon\}$ with $\epsilon=0.05,0.01,0.001$,
        popularised by JP~Morgan's RiskMetrics \ts{13}{14}{1:22:18}. It is standard for cash instruments (stocks, bonds) but,
        as Kempthorne says, for derivatives with nonlinear payoffs ``things get complicated'' \ts{13}{14}{1:24:25}.
  \item \textbf{Conditional VaR} (expected shortfall), $\mathrm{CVaR}_{1-\epsilon}(R_P)=\E[-R_P\mid-R_P\ge\mathrm{VaR}_{1-\epsilon}(R_P)]$, the average
        loss beyond the VaR, which is where VaR is silent; expected to enter bank regulation
        \ts{13}{14}{1:23:22}. The optimisation of CVaR is a linear programme (Rockafellar and Uryasev, 2000, cited on the slide).
\end{itemize}
Definitions with the loss convention, the coherence axioms and the estimation of VaR and ES are in
\cref{ch08:sec-riskmeasures,ch:var}; here only the portfolio-theory consequences are needed.

\begin{proposition}[VaR is not subadditive; variance is not monotone; not stated in class]\label{ch10:prop-var}
Two independent positions each lose $100$ with probability $4\%$ and otherwise $0$. At the $95\%$ level,
$\mathrm{VaR}=0$ for each (the loss occurs with probability $4\%<5\%$), but for the portfolio of both
$\Prob(\text{loss}>0)=1-0.96^2=7.84\%>5\%$, so $\mathrm{VaR}_{95\%}=100>0+0$: diversification \emph{increases} VaR.
The expected shortfalls are $80$ for each and $103.2$ for the portfolio ($\le160$), as CVaR is subadditive. Likewise
the variance violates monotonicity: $R=0$ and $R'\in\{0,1\}$ with probability $\frac12$ each satisfy $R\le R'$ but
$\Var R=0<\Var R'=\frac14$, whereas monotonicity requires $s(R)\ge s(R')$.
\end{proposition}
\noindent (Checked by direct computation: $\mathrm{ES}_{95\%}(\text{portfolio})=(0.16\%\cdot200+4.84\%\cdot100)/5\%=103.2$.)

\begin{coursenote}[Erratum: translation invariance on slide~31 of \file{lecnote14}]
The slide, written for \emph{returns}, defines a coherent measure by: monotonicity ($R_P\le R_{P'}$ w.p.~1 $\Rightarrow s(R_P)\ge s(R_{P'})$),
subadditivity, positive homogeneity, and ``translational invariance $s(R_P+a)\le s(R_P)-a$''. The standard axiom
(Artzner--Delbaen--Eber--Heath 1999) is an \emph{equality}, $s(R_P+a)=s(R_P)-a$: adding a sure gain $a$ lowers the risk by
exactly $a$ (in the loss convention of \cref{ch08:def-coherent}, $\rho(L+a)=\rho(L)+a$). The slide's ``coherent'' verdicts
(variance not, VaR not, CVaR yes) are correct \ts{13}{14}{1:24:25}.
\end{coursenote}

\begin{finance}[Expected loss, worst loss, stop losses, drawdowns]
Asked for the best risk measure, Xia answers the \emph{expected loss} $L$, arguing that gains and losses are
usually not symmetric; one should also know the worst case but cannot size a portfolio on it (one would never invest) \ts{24}{13}{1:14:54}\ts{24}{13}{1:15:58}. Stop losses on trading platforms are set tighter than the
expected loss (cut the position by half after a small loss, another half after the next), which gives the
trader an asymmetric, option-like payoff and forces a distribution of the trades that fits it; and hedge-fund managers,
who take a share of profits but not of the losses, hold ``a free call option'' on the fund
\ts{24}{13}{1:16:58}\ts{24}{13}{1:18:00}\ts{24}{13}{1:19:01}. The loss tolerance of \cref{ch10:sec-problem} is often expressed as a
\emph{maximum drawdown} $\max_{s\le t}\bigl(\max_{u\le s}V_u-V_s\bigr)/\max_{u\le s}V_u$, the largest peak-to-trough fall of
the portfolio value $V$ (not defined in class; see the drawdowns of the concentrated portfolio in \cref{sec:ewport}).
\end{finance}

% ==================================================================
\section{Beyond Markowitz: crowds and power laws (2024)}\label{ch10:sec-crowd}

In 2024 Xia devotes the second half of the lecture to three ideas from his own research (slides marked
``sections 4--7 highlight my research work''): the loss of meaning of volatility as risk (\cref{ch10:sec-kelly}), the
unpredictability of markets because of \emph{crowd behaviour}, and the origin of \emph{power-law} distributions.
They are not tested, and are heuristic, but they answer his ``MPT issue number 2'': how do we predict the future from
the past \ts{24}{13}{54:54}? The scientific method (observe, quantify, extract a pattern, build a model, predict,
repeat and calibrate, control) applies to physics and engineering, but in finance the historical patterns may not
be repeatable because humans adapt and herd \ts{24}{13}{56:01}\ts{24}{13}{57:01}. In 2013 the same point was
made with videos and no formulas: a bridge whose walkers synchronise their steps, and a bank of metronomes on a
moving plate that synchronise \ts{13}{16}{1:05:51}\ts{13}{16}{1:08:02}; in 2024 with a video of birds
\ts{24}{13}{58:03} and the London Millennium Bridge \ts{24}{13}{59:05}. The lesson for portfolio management: an
investor who finds the ``optimal'' strategy finds that everyone else implements it too, and the collective system
is not optimal but fragile; individual managers also have an incentive to herd, since underperforming the peers is
punished more than being wrong with the crowd (career risk)
\ts{13}{16}{1:10:12}\ts{13}{16}{1:16:33}\ts{13}{16}{1:17:37}. Participants change the market they are trying to optimise \ts{13}{16}{1:12:19}.

\subsection{A feedback model of crowd behaviour}
Slide~28 (drawn on slide~34 with heterogeneous agents) describes the interaction loop of $N$ agents, an external force
$E$ and an observation $O$ (the price) \ts{24}{13}{1:00:05}:
\begin{equation}\label{ch10:eq-crowd}
  \dd S_i=C_i\dd E+B_i\dd O+\varepsilon_i,\qquad \dd O=\sum_{i=1}^NA_i\dd S_i ,
\end{equation}
where $S_i$ is the action of agent $i$, $A_i$ its power (influence on the observation), $C_i$ its response to the
external force, $B_i$ its \emph{reactive function} (how strongly it reacts to the observation) and $\varepsilon_i$
noise. In the rational state $B_i$ takes low values $\{B_L\}$; when $\dd O$ moves a lot, agents switch to a reactive
state with high values $\{B_H\}$, and $N_H$ counts them \ts{24}{13}{1:01:08}. The \emph{order parameter}
\[
  R(t)=\frac{\bigl|\sum_i\dd S_i(t)\bigr|}{\sum_i|\dd S_i(t)|}\in[0,1]
\]
is $1$ if all agents act in the same direction (synchronised) and near $0$ if they are not \ts{24}{13}{1:01:08}.
The slides plot $R$ against $N_H/N$ (rising steeply from $\approx0.1$ at $N_H=0$ to $1$ at $N_H=N$) and against the noise
amplitude (falling from $1$ to near $0$: noise desynchronises, as in coupled oscillators)
\ts{24}{13}{1:02:15}, and a simulated ``bubble'': an external force that rises, then falls and oscillates makes the
number of reactive agents rise and fall and the observation follow; if the force stays high forever, all agents end up
reactive and the observation explodes, a \emph{tipping point} when $A\times B_{\max}>1$ (slide~30)
\ts{24}{13}{1:03:21}\ts{24}{13}{1:04:28}.

\begin{proposition}[Mean-field reading of the tipping point; my derivation, not shown in class]\label{ch10:prop-tipping}
Suppose all agents react with the same $B_i=B$. Substituting the first equation of \eqref{ch10:eq-crowd} into the second,
\begin{align*}
  \dd O&=\sum_iA_iC_i\dd E+B\Bigl(\sum_iA_i\Bigr)\dd O+\sum_iA_i\varepsilon_i,\quad\text{so}\\
  \dd O&=\frac{N}{1-\bar A B},\quad N=\sum_iA_iC_i\dd E+\sum_iA_i\varepsilon_i,\quad \bar A=\sum_iA_i .
\end{align*}
The response of the price to any external push or noise is amplified by the factor $1/(1-\bar AB)$, which is finite
and positive iff $\bar AB<1$ and diverges as $\bar AB\to1$: with $\bar A B\ge1$ the loop gain exceeds $1$ and the
feedback is unstable. The switch from $B_L$ to $B_H$ (a higher $B$ when many agents are reactive) moves the system
towards the threshold, which is the tipping point of slide~30.
\end{proposition}
\noindent This is the analogue of the gain of a linear amplifier with positive feedback (Barkhausen criterion), or of the
critical coupling of a mean-field ferromagnet, $\chi=1/(1-J\bar A)$, the susceptibility diverging at the critical
temperature. The lecture presents the mechanism qualitatively; the numerical scenarios on the slides are simulations
whose parameters are not given.

\subsection{Power laws}
Why are wealth, city sizes, venture-capital returns or viral posts so unequal, unlike heights?
\ts{24}{13}{1:05:33}\ts{24}{13}{1:06:33}. The $80$--$20$ rule, the Matthew effect and ``winner takes all'' describe distributions
that are \emph{scale-free}: the top $20\%$ of the people own $80\%$ of the wealth, and within them again the top $20\%$
own $80\%$, and so on. In symbols $p(bx)=g(b)p(x)$ for all $b>0$ holds (with smooth $p$) iff $p$ is a power law, $p(x)=Cx^{-a}$
(with $g(b)=b^{-a}$; the ``only if'' is the standard functional-equation argument, $p(bx)=g(b)p(x)$ forces $g(bb')=g(b)g(b')$) \ts{24}{13}{1:05:33}.
The three usual descriptions are related as follows (slide~33) \ts{24}{13}{1:09:39}\ts{24}{13}{1:10:39}:
\begin{align*}
  \text{Pareto (tail):}&\quad \Prob[X>x]\sim x^{-k}, & \text{Zipf (rank $r$):}&\quad \E[x_{(r)}]\sim r^{-b},\\
  \text{density:}&\quad f(x)=Cx^{-a}, & a&=1+k=1+\frac1b .
\end{align*}
(Differentiating the tail gives $a=1+k$; the $r$-th largest of $N$ satisfies $r/N=\Prob[X>x_{(r)}]$, so $x_{(r)}\sim r^{-1/k}$ and $b=1/k$.)
A Gaussian has no such structure; heights are Gaussian because they are set once, whereas in social
phenomena ``the rich get richer'': agents have \emph{different powers} $A(i,t)$ and gain power by betting in the
right direction of the observation, so the power of the crowd concentrates in a few ``super agents'' (governments and central banks, the Fed being the example named in class; in the old days, some hedge funds) whose actions dominate the system
\ts{24}{13}{1:08:38}\ts{24}{13}{1:10:39}\ts{24}{13}{1:12:49}. He also mentions entropy and stability of such systems as a topic
``beyond the scope''.

\begin{example}[The $80$--$20$ rule and the Pareto exponent]\label{ch10:ex-pareto}
If wealth has the tail $\Prob[X>x]=(x_m/x)^k$ with $k>1$, the top fraction $q$ of the population holds the fraction
$q^{1-1/k}$ of the total wealth (integrate $x$ over the tail). The $80$--$20$ rule, $0.2^{1-1/k}=0.8$, gives
$k=\ln5/\ln4=1.161$, $a=2.16$, $b=0.861$; and the top $20\%$ of the top $20\%$ hold $80\%$ of the $80\%$, i.e.\ $64\%$ of all
wealth (top $4\%$), the self-similarity of the lecture. Note $k$ close to $1$: the mean is barely finite
(for $k\le1$ it is infinite).
\end{example}

\begin{finance}[What the crowd material means for practice]
Xia's conclusions of the 2024 lecture (slide~35) \ts{24}{13}{1:11:45}: portfolio management is about \emph{sizing}; clarify
objectives and know the loss tolerance; diversification helps but requires rebalancing; volatility is not risk and the
Sharpe ratio captures neither skewness nor sizing (use the gain--loss ratio); capital market assumptions are
not reliable; crowd behaviour can push markets to extremes; watch the powerful super agents; understand the key drivers of the
portfolio. The 2013 version ends similarly: ``solving the problem'' is not done, since \emph{when you participate
in the market you change the market}, and quantitative finance ``is not like solving physics problems''
\ts{13}{16}{1:11:12}\ts{13}{16}{1:12:19}. A questioner asked whether this is game theory; Xia agreed that it has
a lot to do with it, but that the markets have no clearly defined rules \ts{13}{16}{1:13:22}. Whether a
computer can beat a human in investment, he ``cannot confidently say it will not happen''; a colleague added that
managers also manage people \ts{13}{16}{1:18:37}.
\end{finance}


% ==================================================================
\section{Case study: simulation and US sector ETFs (2013 Case Study~6)}\label{ch10:sec-case}

\begin{casestudy}[\file{CaseStudy6.pdf}, \file{Smltn_TwoAst.pdf}, \file{ETF_prid*.pdf}]
\textbf{Goal.} Kempthorne's live demonstration of the theory in R \ts{13}{14}{8:32}\ts{13}{14}{1:04:54}:
\begin{itemize}
  \item the feasible set of two
      assets for several correlations;
  \item mean--variance optimal portfolios of the nine US sector ETFs, with a riskless
      asset, with the constraint that no asset exceeds $30\%$ or $15\%$ of the capital.
\end{itemize}

\textbf{Part 1: two-asset simulation.} Two Gaussian assets with $\mu=(0.15,0.25)$, $\sigma=(0.20,0.30)$ (see the erratum
below), 500 weekly returns simulated for $\rho=-0.8,-0.4,0,0.4,0.8$. Each page of \file{Smltn_TwoAst} has four
panels: cumulative returns of the two assets (blue: asset~2, green: asset~1); the feasible set of
\cref{ch10:sec-two}; the scatter plot of the weekly returns; and the same cumulative returns with the cumulative return
of the minimum-variance portfolio (cyan) \ts{13}{14}{8:32}\ts{13}{14}{10:44}. In the PDF as distributed only the two asset
points are drawn in the second panel and the scatter panel is blank, so the curves of \cref{ch10:fig-two} were recomputed.
The sample statistics printed above the panels are:
\begin{center}
\small
\begin{tabular}{@{}lrrrrr@{}}
\toprule
$\rho$ & $-0.8$ & $-0.4$ & $0$ & $0.4$ & $0.8$\\
\midrule
sample means & $0.115,\ 0.341$ & $0.144,\ 0.343$ & $0.183,\ 0.334$ & $0.210,\ 0.321$ & $0.085,\ 0.177$\\
sample vols & $0.191,\ 0.291$ & $0.194,\ 0.291$ & $0.200,\ 0.289$ & $0.200,\ 0.287$ & $0.196,\ 0.287$\\
sample corr.\ & $-0.784$ & $-0.359$ & $0.031$ & $0.395$ & $0.791$\\
\bottomrule
\end{tabular}
\end{center}
Sample volatilities and correlations are close to the population values; the sample means are off by about one
standard error (\cref{ch10:sec-est}); for $\rho=0.8$ both are well below the population values ($0.085$ and $0.177$ against $0.15$ and $0.25$), an unlucky
draw, and the cumulative-return curves of the last panel look very different from the others. Lessons observed live: with
$\rho\le0.4$ the minimum-variance portfolio has a return and a volatility better than asset~1 alone; with
$\rho=-1$ the variance would be zero \ts{13}{14}{11:47}\ts{13}{14}{17:06}; every portfolio to the left of the minimum-variance point is
suboptimal \ts{13}{14}{14:57}.

\textbf{Part 2: nine sector ETFs.} The SPDR sector funds XLB (materials), XLV (health care), XLP (consumer staples),
XLY (consumer discretionary), XLE (energy), XLF (financials), XLI (industrials), XLK (technology) and XLU (utilities), for two
periods: 2009--2013 (``periodA'', up to a week before the lecture; $251$ observations on the time axis, i.e.\ weekly
data over about $4.8$ years, an inference: the frequency is not stated) and 2003--2006 (``periodB'', $209$ observations).
\cref{ch10:tab-etf} gives the annualised return and volatility of each. The correlations are high, with the
average pairwise correlation $0.74$ in 2009--2013 (range $0.57$--$0.93$) and $0.55$ in 2003--2006 (range
$0.16$--$0.85$): there is much less diversification available in the aftermath of the crisis.

\textbf{Method.} For a grid of $101$ target returns (``target return index'' in the plots) the R script minimises the variance
of a portfolio of the nine ETFs and the riskless asset (return $\approx0$: the frontier of plot~4 starts at the origin),
subject to full investment, $w_j\ge0$ and $w_j\le30\%$ or $15\%$. It prints the allocations for three target volatilities and
draws six plots: \emph{1}~risk--return scatter with the riskless asset; \emph{2}~optimal allocations versus target
return (lines); \emph{3}~the same as a stacked bar chart (dark red at the top is cash); \emph{4}~the efficient frontier over
the ETFs; \emph{5}~the frontier with the target portfolio for a $10\%$ volatility (circle); \emph{6}~cumulative returns of the
ETFs with that optimal portfolio in bold blue \ts{13}{14}{1:04:54}--\ts{13}{14}{1:11:11}.
\end{casestudy}

\begin{table}[ht]
\centering
\begin{tabular}{@{}llrrrr@{}}
\toprule
& & \multicolumn{2}{c}{2009--2013} & \multicolumn{2}{c}{2003--2006}\\
\cmidrule(lr){3-4}\cmidrule(lr){5-6}
ETF & sector & return & vol. & return & vol.\\
\midrule
XLB & materials & 0.16 & 0.24 & 0.16 & 0.18\\
XLV & health care & 0.17 & 0.15 & 0.07 & 0.12\\
XLP & consumer staples & 0.15 & 0.12 & 0.08 & 0.09\\
XLY & consumer discretionary & 0.24 & 0.21 & 0.14 & 0.14\\
XLE & energy & 0.15 & 0.25 & 0.26 & 0.20\\
XLF & financials & 0.13 & 0.33 & 0.15 & 0.13\\
XLI & industrials & 0.18 & 0.23 & 0.15 & 0.14\\
XLK & technology & 0.18 & 0.19 & 0.12 & 0.18\\
XLU & utilities & 0.10 & 0.16 & 0.19 & 0.13\\
\bottomrule
\end{tabular}
\caption{Annualised return and volatility of the nine sector ETFs printed in \file{CaseStudy6}. (The correlation
matrices are printed there as well; the computations of \cref{ch10:ex-etf-unc} and of this section use them.)}\label{ch10:tab-etf}
\end{table}

\paragraph{Results (2009--2013).} With a $30\%$ cap and a target volatility of $9.9\%$ the printed optimum is XLV $15.7\%$,
XLP $30\%$, XLY $22.6\%$, cash $31.7\%$, for a return of $12.4\%$; for a target volatility of $15.3\%$ it is fully
invested: XLV $30\%$, XLP $22.7\%$, XLY $30\%$, XLK $17.3\%$, for a return of $18.7\%$; for $0.9\%$ volatility it
holds $93.5\%$ in cash. The comments of the case study, confirmed by the plots \ts{13}{14}{1:07:59}--\ts{13}{14}{1:11:11}:
\begin{itemize}
  \item As the target return rises from zero only XLY, XLP and XLV enter, in \emph{fixed proportions} (the de-levered
        long-only tangency portfolio, \cref{ch10:prop-constr}; the unconstrained one of \cref{ch10:sec-est} is very
        different); allocations grow linearly (plot~2).
  \item The cap binds first for consumer staples; XLY and XLV are increased; when all three are at $30\%$ technology
        (higher return) enters and the allocations to XLY and XLP fall to make room. No allocation is ever given to financials.
  \item In the efficient frontier of plot~4, every ETF except consumer discretionary is dominated by an optimal
        portfolio, and the frontier shows the CML-like straight segment up to the first kink.
  \item With a $15\%$ cap the constraints bind sooner: XLV, XLP, XLY, XLK, then XLU and XLI, and finally XLB; the frontier is
        \emph{below} the $30\%$ one (\cref{ch10:prop-constr}(i)), because to reach a higher return one has to
        hold assets with a lower return per unit of risk \ts{13}{14}{1:12:14}\ts{13}{14}{1:14:24}.
  \item In 2003--2006 the optimal portfolios are completely different: utilities, energy, financials, consumer
        staples and industrials, XLU being the first to hit the cap, and technology never enters ($15\%$ cap:
        XLP, XLY, XLE, XLF, XLI, XLU at the cap, XLB in addition at the top). The ranking of sectors in one period says
        little about the next.
\end{itemize}

\begin{figure}[ht]
\centering
\begin{tikzpicture}
\begin{axis}[width=0.9\linewidth,height=7.4cm,xlabel={volatility},ylabel={expected return},
  xmin=0.05,xmax=0.36,ymin=0.04,ymax=0.32,xtick={0.1,0.15,0.2,0.25,0.3,0.35},scaled ticks=false,/pgf/number format/fixed,grid=major,grid style={black!10},
  legend style={font=\footnotesize,at={(0.5,1.03)},anchor=south,legend columns=3},
  tick label style={font=\small},label style={font=\small}]
  \addplot[thick,thmcol,dashed] coordinates {(0.05,0.0927) (0.16,0.297)};
  \addplot[very thick,defcol] coordinates {(0.005,0.0065) (0.02,0.0258) (0.04,0.0516) (0.06,0.0774) (0.08,0.1030) (0.10,0.1281) (0.12,0.1530) (0.13,0.1654) (0.14,0.1767) (0.15,0.1867) (0.16,0.1911)};
  \addplot[very thick,intcol] coordinates {(0.005,0.0065) (0.02,0.0258) (0.04,0.0515) (0.06,0.0765) (0.08,0.0991) (0.10,0.1183) (0.12,0.1350) (0.13,0.1433) (0.14,0.1514) (0.15,0.1590) (0.16,0.1664)};
  \legend{unconstrained CML (Sharpe $1.85$),cap $30\%$,cap $15\%$}
  \addplot[only marks,mark=*,mark size=1.8pt,black,forget plot] coordinates {(0.24,0.16) (0.15,0.17) (0.12,0.15) (0.21,0.24) (0.25,0.15) (0.33,0.13) (0.23,0.18) (0.19,0.18) (0.16,0.10)};
  \node[font=\scriptsize,anchor=north east] at (axis cs:0.236,0.157) {XLB};
  \node[font=\scriptsize,anchor=south east,fill=white,fill opacity=0.8,inner sep=1pt] at (axis cs:0.146,0.173) {XLV};
  \node[font=\scriptsize,anchor=north west] at (axis cs:0.123,0.146) {XLP};
  \node[font=\scriptsize,anchor=south] at (axis cs:0.21,0.243) {XLY};
  \node[font=\scriptsize,anchor=north] at (axis cs:0.25,0.147) {XLE};
  \node[font=\scriptsize,anchor=north] at (axis cs:0.33,0.127) {XLF};
  \node[font=\scriptsize,anchor=south west] at (axis cs:0.226,0.185) {XLI};
  \node[font=\scriptsize,anchor=south] at (axis cs:0.19,0.183) {XLK};
  \node[font=\scriptsize,anchor=north] at (axis cs:0.16,0.097) {XLU};
\end{axis}
\end{tikzpicture}
\caption{Nine sector ETFs, 2009--2013 (inputs of \cref{ch10:tab-etf} and the printed correlations; the risk-free rate is
taken to be $0$). Dots: individual ETFs. Blue and green: efficient frontier with cash and long-only positions capped at $30\%$
and $15\%$ (plots~4 of \file{ETF_pridA_30} and \file{ETF_pridA_15}, recomputed here). The dashed line is the
CML of the \emph{unconstrained} problem, with a Sharpe ratio of $1.85$, reached only with large short positions
(\cref{ch10:sec-est}).}\label{ch10:fig-etf}
\end{figure}

\begin{coursenote}[Independent check of the case study]
The R code is not available, so the printed allocations were re-derived from the printed means, volatilities and
correlations (all rounded to two decimals) by maximising the return at a given volatility, with $r_0=0$,
$w_j\in[0,c]$ and full investment or less, with a general-purpose solver. At a target volatility of $15.3\%$ ($30\%$ cap) it gives XLV $30\%$, XLP $22.9\%$, XLY $30\%$,
XLK $17.1\%$ and a return of $18.8\%$ (printed: $30$, $22.7$, $30$, $17.3$, $18.7$); at $9.9\%$ it gives XLV $21.8\%$, XLP $30\%$, XLY $18.7\%$,
cash $29.5\%$ and $12.7\%$ (printed: $15.7$, $30$, $22.6$, $31.7$, $12.4$). The sets of assets, the order in which the caps bind
and the frontiers agree with the printed ones; the few percentage points of difference in XLV and XLY
come from the rounding of the inputs (a half-point change in a mean moves weights by several points), a small illustration of the sensitivity of \cref{ch10:sec-est}. For the $15\%$ cap at $9.9\%$ volatility the recomputation gives XLV, XLP, XLY, XLK at $15\%$,
XLU $3.6\%$, XLI $1.6\%$, cash $34.8\%$ (printed: XLU $6.9\%$, cash $33.1\%$, return $11.7\%$ in both).
\end{coursenote}

\begin{finance}[Sector market-neutral models and diversification]
In the second demonstration of 2013 L14 Kempthorne replaces the ETFs by hedge-fund style strategies: multi-factor pricing
models trading long/short, market-neutral, within each of the nine sectors, over five years
\ts{13}{14}{1:15:29}\ts{13}{14}{1:16:33}. Their total returns are modest ($20$--$60\%$ over five years) but they are
\emph{little correlated} with each other, so the optimiser gets a large benefit from combining them; the model for the utilities sector
receives the largest weight, then industrials and energy, and a target volatility of $10\%$ is reached \ts{13}{14}{1:17:35}\ts{13}{14}{1:18:36}.
The data behind the plots are not in the files, so only this qualitative description is given. The lesson is the one
of \cref{ch10:sec-two}: low correlation makes the frontier bend to the left, and diversification is worth most for
uncorrelated bets.
\end{finance}

% ==================================================================
\section{Notes on the sources, differences between the years and errata}\label{ch10:sec-notes}

\begin{coursenote}[2013 versus 2024]
The 2013 lectures are split: Kempthorne gives the mathematics (his lecture has a full derivation, a case study,
utility theory, constraints and risk measures), Xia the practice. The 2024 lecture has only Xia, and it is a
different lecture from his 2013 one. It keeps the special cases of two assets (slides 14--20; Xia says he skips the mathematics that he covered in earlier years \ts{24}{13}{40:00}) and the rebalancing example, and it adds: the endowment model, MPT
issue~1 (volatility as risk) with the gain--loss ratio, MPT issue~2 (predicting the future) with the crowd model, and
power laws. It drops the Kelly formula, the beta discussion and the Millennium Bridge \emph{and} metronome videos of 2013 (2024 shows
a video of birds and mentions the bridge). The 2024 slides point to the 2013 video of Xia's lecture as
``one of the top five most viewed courses'' of OpenCourseWare \ts{24}{13}{1:06}. What Kempthorne derives (Markowitz problem, tangency
portfolio, Tobin, utility, constraints, risk measures) is not covered in 2024 by anyone: the derivations of this chapter
follow \file{lecnote14}.
\end{coursenote}

\begin{coursenote}[Erratum: parameters of the two-asset simulation]
Kempthorne says at the start of the demonstration that asset~1 has mean $15\%$ and volatility $25\%$, asset~2 mean $20\%$ and
volatility $30\%$ \ts{13}{14}{5:21}, and page~2 of \file{CaseStudy6} prints the same numbers (with typos:
``$E(R_1)=0.20=\alpha_2$'' and ``$\sqrt{\Var R_1}=0.30=\sigma_2$'' should refer to $R_2$, and ``$\sigma_w^2=\Var(R_2)$'' should be
$\Var(R_w)$). But the plots of \file{Smltn_TwoAst} are titled ``mean returns $0.15,0.25$, asset volatilities $0.20,0.30$'',
Kempthorne himself reads ``mean returns $15\%$, $25\%$, volatility $20$ to $30$'' when he shows them \ts{13}{14}{10:44}, and the
sample statistics ($0.19$--$0.20$ and $0.29$ for the volatilities, cumulative return $\approx3.3$ for asset~2, i.e.\ about $0.34$ per year over $9.6$ years) fit
these values, not the printed ones. This chapter uses the values of the plots.
\end{coursenote}

\begin{coursenote}[Errata and remarks on \file{CaseStudy6}]
\begin{itemize}
  \item Section~2.3 refers to ``Plot~4 in the two files \file{ETF_S_1_periodA_30.pdf} and \file{ETF_S_1_periodA_30.pdf}'': the second
      should be the $15\%$ file (\file{ETF_pridA_15}); the case study cites the file names of the original R output
      (\texttt{ETFS\_1\_periodA\_30}), which differ from the OCW names \file{ETF_pridA_30}.
  \item Section~3.3 (2003--2006, $15\%$ cap) says both ``technology
      never enters the optimal allocation'' and ``the allocations to XLK increase until its limit is reached'': the second
      bullet is copied from Section~2.3; the printed table has XLK at $0.000$ throughout.
  \item The bars of plot~3 use very
      similar colours for consumer discretionary and energy (both green) and the line plots repeat colours (XLB and XLU are both
      black in the first plot of each file), so the ETFs are identified through the tables.
\end{itemize}
\end{coursenote}

\begin{coursenote}[Slides and lectures: small corrections]
\begin{itemize}
  \item \file{lecnote14}, slide~19 writes $\E[u(W)]\propto\E[W-\frac12\lambda(W-w_*)^2]\approx\E[W]-\frac12\lambda\Var W$: the
      $\propto$ hides an additive constant and the positive factor $u'(w_*)$; the content is that of \cref{ch10:prop-mv-utility}(i).
  \item Slide~31, translation invariance: see the erratum in \cref{ch10:sec-risk}.
  \item Slide~24 of \file{lec13} asserts the ratio $(G-L)/(G+L)$ is ``the same as
      the Kelly criterion for binary bets'': true only for even-money bets (\cref{ch10:prop-gl}(i)).
  \item In 2013 L16 Xia says that after year~1 the winning
      asset ``becomes 75'' when rebalancing: the correct amounts are $62.5$ in each asset (\cref{ch10:sec-rebal}); the 2024 slide has it right.
  \item The transcript of 2013 L14 says Fama ``was just awarded, what, two weeks ago'': the 2013 economics prize
      (Fama, Hansen, Shiller) was announced on 14 October 2013, ten days before the lecture of 24 October (the case-study date).
  \item In 2013 L16 two remarks made ``from the audience'' (about estimating returns versus volatility, and about the Kelly criterion, Shannon and Samuelson)
      are by a co-lecturer, in all likelihood Kempthorne; the transcript does not identify the speaker.
  \item \file{lecnote16} contains only the outline (five items: portfolio construction, portfolio theory in special cases, risk parity, limitations, what we learned): Xia used a single slide and the blackboard, so 2013 L16
      has no formulas apart from those reconstructed here from the transcript and from the 2024 slides.
\end{itemize}
\end{coursenote}

% ==================================================================
\begin{keyformulas}
\begin{align*}
&\text{Portfolio:} && \mu_{\bm w}=\bm w\T\bm\mu,\quad \sigma_{\bm w}^2=\bm w\T\Sigma\bm w,\quad \bm w\T\bm 1=1\\
&\text{Two assets:} && \sigma_w^2=(1-w)^2\sigma_1^2+w^2\sigma_2^2+2w(1-w)\rho\sigma_1\sigma_2\\
& && w^*=\tfrac{\sigma_1^2-\rho\sigma_1\sigma_2}{\sigma_1^2+\sigma_2^2-2\rho\sigma_1\sigma_2},\quad
  \sigma^{*2}=\tfrac{\sigma_1^2\sigma_2^2(1-\rho^2)}{\sigma_1^2+\sigma_2^2-2\rho\sigma_1\sigma_2}\\
&\text{Constants:} && a=\bm\mu\T\Sigma^{-1}\bm\mu,\ b=\bm\mu\T\Sigma^{-1}\bm 1,\ c=\bm 1\T\Sigma^{-1}\bm 1,\ \Delta=ac-b^2\\
&\text{Optimal weights:} && \bm w_0=\lambda_1\Sigma^{-1}\bm\mu+\lambda_2\Sigma^{-1}\bm 1,\quad
  (\lambda_1,\lambda_2)=\tfrac1\Delta(c\mu_0-b,\ a-b\mu_0)\\
&\text{Frontier:} && \sigma_0^2=\tfrac{c\mu_0^2-2b\mu_0+a}\Delta=\tfrac1c+\tfrac c\Delta(\mu_0-\tfrac bc)^2\\
&\text{Min.\ variance:} && \bm w_g=\tfrac{\Sigma^{-1}\bm 1}c,\quad\sigma_g^2=\tfrac1c,\quad \mu_g=\tfrac bc\\
&\text{Risk aversion:} && \bm w(\lambda)=\bm w_g+\tfrac1\lambda\Sigma^{-1}(\bm\mu-\tfrac bc\bm 1)\\
&\text{With riskless asset:} && \bm w_0=\tfrac{\mu_0-r_0}{d^2}\Sigma^{-1}\bm\eta,\quad \bm\eta=\bm\mu-r_0\bm 1,\quad d^2=\bm\eta\T\Sigma^{-1}\bm\eta\\
&\text{Tangency portfolio:} && \bm w_M=\tfrac{\Sigma^{-1}\bm\eta}{\bm 1\T\Sigma^{-1}\bm\eta},\quad \mu_M-r_0=\tfrac{d^2}{b-r_0c},\quad
  \sigma_M=\tfrac d{b-r_0c}\\
&\text{Tobin, CML:} && m_{\bm w}=\tfrac{\mu_0-r_0}{\mu_M-r_0},\quad \mu_P=r_0+\tfrac{\mu_M-r_0}{\sigma_M}\sigma_P\\
&\text{Sharpe ratio:} && \mathrm{SR}=\tfrac{\mu_P-r_0}{\sigma_P}\le d\ \text{(max.\ at }\bm w_M),\ \text{leverage-invariant}\\
&\text{SML:} && \mu_i-r_0=\beta_i(\mu_M-r_0),\quad \beta_i=\tfrac{\Cov(R_i,R_M)}{\sigma_M^2}=\rho_{iM}\tfrac{\sigma_i}{\sigma_M}\\
&\text{Equal weights:} && \Var R_{\rm ew}=\tfrac{\bar\sigma^2}n+(1-\tfrac1n)\bar c\\
&\text{Constant mix:} && g(\bm w)=\bm w\T\bm\mu-\tfrac12\bm w\T\Sigma\bm w\ \ (\text{log growth})\\
&\text{Risk contribution:} && \mathrm{RC}_i=\tfrac{w_i(\Sigma\bm w)_i}{\sigma_P},\ \textstyle\sum_i\mathrm{RC}_i=\sigma_P;\ \text{2 assets: }w_i\propto1/\sigma_i\\
&\text{Utility:} && \lambda_A=-\tfrac{u''}{u'},\ \lambda_R=-\tfrac{Wu''}{u'};\quad \E[u]\approx u(w_*)+\tfrac12u''\Var W;\\
& && \text{CARA+normal: CE}=\mu_W-\tfrac\lambda2\sigma_W^2\\
&\text{Kelly:} && f^*=\tfrac pl-\tfrac qg;\ \ \text{Gaussian: }f^*=\tfrac{\mu}{\sigma^2}\\
&\text{Gain--loss ratio:} && \tfrac{G-L}{G+L}=\tfrac{\E[R]}{\E[|R|]}\in[-1,1],\ G=\E[R^+],\ L=\E[R^-]\\
&\text{Constraints:} && \mathrm{TE}^2=(\bm w-\bm w_B)\T\Sigma(\bm w-\bm w_B),\ \ \textstyle\sum_i\beta_{ik}w_i=0\\
&\text{Estimation:} && \mathrm{se}(\hat\mu)=\tfrac\sigma{\sqrt T},\quad \tfrac{\mathrm{se}(\hat\sigma)}\sigma\approx\tfrac1{\sqrt{2n}}\\
&\text{Crowd loop:} && \dd O=\tfrac{\sum A_iC_i\dd E+\sum A_i\varepsilon_i}{1-\bar AB}\ \ (\text{mean field})\\
&\text{Power law:} && \Prob[X>x]\sim x^{-k},\ f\sim x^{-a},\ a=1+k=1+\tfrac1b
\end{align*}
\end{keyformulas}

% ==================================================================
\section*{Exercises}
\addcontentsline{toc}{section}{Exercises}

\subsection*{From 2013 Problem Set 6 (Problems 3--5; Problems 1--2 are in \cref{ch:timeseries})}
In all three problems the portfolio $\bm w$ is fully invested, $\sum_jw_j=1$, with no short sales, $w_j\ge0$.

\begin{exercise}[PS6 (2013), Problem 3: two risky assets]\label{ch10:ex-ps6-3}
Two assets with returns $\bm R=(R_1,R_2)\T$, means $(\alpha_1,\alpha_2)$, covariance
$\Sigma=\begin{psmallmatrix}\sigma_1^2&\rho\sigma_1\sigma_2\\ \rho\sigma_1\sigma_2&\sigma_2^2\end{psmallmatrix}$ and a portfolio
$R_{\bm w}=w_1R_1+w_2R_2$.
\begin{enumerate}[(a)]
  \item Prove that $\Var(R_{\bm w})\le\max(\sigma_1^2,\sigma_2^2)$ for every such portfolio.
  \item For $\sigma_1=\sigma_2$ and $\rho=0$ find the minimum-variance portfolio $\bm w^*$ and its variance.
  \item For $\sigma_1=\sigma_2$ and arbitrary $\rho$ find $\bm w^*$ and $\Var(R_{\bm w^*})$, and graph the minimal variance as a function of
      $\rho\in[-1,1]$.
\end{enumerate}
\end{exercise}

\begin{exercise}[PS6 (2013), Problem 4: $m>2$ assets, equal volatilities]\label{ch10:ex-ps6-4}
Now $m>2$ assets with covariance matrix $\Sigma=(\rho_{ij}\sigma_i\sigma_j)$.
\begin{enumerate}[(a)]
  \item Prove that $\Var(R_{\bm w})\le\max_j\sigma_j^2$ for every long-only fully invested portfolio.
  \item If $\sigma_1=\dots=\sigma_m$ and all $\rho_{ij}=0$ ($i\neq j$), find the minimum-variance portfolio, express its variance as a function
      of $m$ and find its limit as $m\to\infty$.
  \item If $\sigma_1=\dots=\sigma_m$ and $\rho_{ij}\equiv\rho$ for all $i\neq j$ (\emph{equicorrelation}): first show that
      $\rho\ge-1/(m-1)$ (hint: $\bm 1$ is an eigenvector of $\Sigma$; compute its eigenvalue and impose that $\Sigma$ be positive
      semi-definite). Then find the minimum-variance portfolio and its variance, graph the variance as a function of $\rho\in[-\frac1{m-1},1]$,
      find the limit as $m\to\infty$ and compare it with (b), commenting on the ability to diversify away risk by adding equicorrelated assets.
\end{enumerate}
\end{exercise}

\begin{exercise}[PS6 (2013), Problem 5: unequal variances and the general case]\label{ch10:ex-ps6-5}
\leavevmode
\begin{enumerate}[(a)]
  \item For $m=2$, $\sigma_1\neq\sigma_2$ and $\rho=0$, find the minimum-variance portfolio and its variance.
  \item For $m>2$ and $\Sigma=\diag(\sigma_1^2,\dots,\sigma_m^2)$ (arbitrary variances, zero correlations) find the minimum-variance
      portfolio and its variance, and express the variance in terms of $m$ and $\tilde\sigma^2=\bigl(\frac1m\sum_j\sigma_j^{-2}\bigr)^{-1}$.
  \item For $m>2$ and an arbitrary positive definite $\Sigma$, find the minimum-variance portfolio and its variance.
\end{enumerate}
\end{exercise}

\begin{exercise}[2024 Investment Game and the gain--loss ratio (ungraded)]\label{ch10:ex-game}
The 2024 lecture (slide~36) repeats the game of \cref{ex:game}: $\$10{,}000$ in one stock or ETF from the close of 13~September to the
close of 15~November, one optional switch on 11~October, and the final calculation of $G$, $L$ and $(G-L)/(G+L)$ from the daily
profits and losses. Carry it out on a two-month window of your choice (for example for the sector ETFs of \cref{ch10:tab-etf}), compute the
Sharpe ratio of the same daily P\&L, and compare the ranking of two investments by the two measures; see \cref{ch10:prop-gl}.
\end{exercise}

\subsection*{Extra exercises (not from the course)}

\begin{exercise}[not from the course: risk parity]
Two assets have volatilities $16\%$ and $6\%$ and correlation $0.3$.
\begin{enumerate}[(a)]
  \item Find the risk-parity weights and check that the risk
      contributions of \cref{ch10:def-rc} are equal.
  \item What fraction of the risk of the $60/40$ portfolio comes from the first asset?
  \item By what factor must the
      risk-parity portfolio be levered to have the volatility of the $60/40$ portfolio?
\end{enumerate}
\end{exercise}

\begin{exercise}[not from the course: the frontier]
\leavevmode
\begin{enumerate}[(a)]
  \item Verify that $\bm w_g$ of \cref{ch10:cor-gmv} minimises $\bm w\T\Sigma\bm w$ under $\bm w\T\bm 1=1$ and that the frontier
      \eqref{ch10:eq-frontier} has asymptotes $\mu=b/c\pm\sqrt{\Delta/c}\,\sigma$.
  \item Show that the difference of two efficient portfolios is proportional to $\bm v$ of
      \eqref{ch10:eq-twofund} and deduce the two-fund theorem.
  \item Show that, if $r_0<\mu_g$, the tangency portfolio maximises the
      Sharpe ratio over the fully invested portfolios, and find the tangency point of the line through $(0,r_0)$ with the hyperbola.
\end{enumerate}
\end{exercise}

\begin{exercise}[not from the course: CARA utility]
An investor with $u(W)=1-e^{-AW}$ and initial wealth $W_0$ can invest in $m$ assets with normal returns and in a riskless asset. Show that the
optimal amounts invested in the risky assets are $A^{-1}\Sigma^{-1}\bm\eta$ (in dollars) and identify the fraction of the wealth held in the
tangency portfolio.
\end{exercise}

\begin{exercise}[not from the course: Kelly]
A bet wins $g=2$ per unit staked with probability $p=0.6$ and loses $l=1$ otherwise.
\begin{enumerate}[(a)]
  \item Compute the Kelly fraction and the
      optimal growth rate $G(f^*)$.
  \item What growth rate does the half-Kelly fraction $f^*/2$ give, as a percentage of $G(f^*)$?
  \item Show that for small edges
      half-Kelly gives $3/4$ of the optimal growth rate.
\end{enumerate}
\end{exercise}

\begin{exercise}[not from the course: gain--loss ratio]
Investment~A pays $+3$ with probability $0.4$ and $-1$ with probability $0.6$. Investment~B pays $+1$ or $-1$
with probability $0.5$ each plus a fixed $0.3$. Compute for each the mean, the standard deviation, the Sharpe ratio, and
$(G-L)/(G+L)$; comment on the difference in rankings and on the Kelly fractions.
\end{exercise}

\begin{exercise}[not from the course: VaR and CVaR]
Modify \cref{ch10:prop-var} to three independent positions with default probability $3\%$ and loss $100$: compute the
$95\%$ VaR and expected shortfall of one position and of the portfolio of three, and check subadditivity of ES.
\end{exercise}

\begin{exercise}[not from the course: rebalancing]
Two independent assets have geometric Brownian motions with $\sigma_1=30\%$, $\sigma_2=10\%$ and the same $\mu$. Using
\cref{ch10:prop-growth}, find the constant-mix weight that maximises the log growth rate and the value of the ``excess growth'' at
that weight. Compare with equal weights.
\end{exercise}

\begin{exercise}[not from the course: crowd model and power laws]
\leavevmode
\begin{enumerate}[(a)]
  \item In \cref{ch10:prop-tipping}, take $N$ agents with $A_i=A/N$ and $B_i\in\{B_L,B_H\}$, a fraction $x$ in the reactive state; write the loop
      gain and find the critical $x$ for the tipping point.
  \item For a Pareto tail $\Prob[X>x]=(x_m/x)^k$, $k>1$, prove that the top fraction $q$
      holds $q^{1-1/k}$ of the total, and compute $k$ for a ``$90$--$10$'' rule.
\end{enumerate}
\end{exercise}
